% ============================================================
% PAPER 5 of 11 -- exact_belief_computation
% Assembled 2026-09-29 by content-and-merits partition.
% Title: Exact Belief Computation on the Observation Ladder
% Source: arena agent 1/paper rewrites/latex/paper2_exact_belief_computation_v13.tex
% Partition note: Single-source. Short self-contained result: Hamming classification and the antichain census.
%% STATUS: structurally assembled as a NEW version. Editorial integration
%%         (de-duplication of shared framework, sibling cross-citations)
%%         is NOT yet done. See PAPERS_MANIFEST.md.
%% ============================================================

%% v12 (2026-09-29): (a) SCOPE DECISION RESOLVED -- expand, do not fold into paper 3.
%%   Paper 3 CITES this paper as "the exact-computation companion"; folding would join
%%   two different mathematical cores and break that citation.
%% (b) 'ebc' identification CONFIRMED, moderate->high: 16 Lean modules named EBC_*
%%   (Bands/Classification/Deadline/Dynamics/ExactBelief/Hamming/Ladder/Pairs) in
%%   namespace Formalizations.EBC map onto this paper's own sections.
%% (c) 'Scale II' relation RESOLVED: the parent theory (belief-state safety value
%%   V_k(b)) is paper2_belief_state_v2 -- paper 2 of this set in its later form,
%%   32.6%% of its sentences appear verbatim in paper02. Stated in the text.
%% (d) Added subsection priorart: concedes Sperner/Dedekind/Kleitman and Hamming
%%   adjacency as classical, states this paper proves NO new theorem about
%%   antichains, narrows the claim to four items. Added subsection mech (16 EBC
%%   modules, 134 declarations, 0 sorry) with the honest note that the census
%%   enumeration itself is NOT formalized.
%% (e) Markups are \subsection to match this paper (it uses no \section).
%% PRIOR ART AND SCOPE REQUIRED -- bar item, inserted 2026-09-29.
%% This is the shortest paper in the family (~3.6k words) with no related-work
%% position, and it is the furthest from "substantial methodological development".
%% Two things must happen before it can stand:
%%   (i) PRIOR ART: exact and rational-arithmetic computation of viability
%%       objects -- Saint-Pierre's kernel algorithms and their successors -- and
%%       the combinatorial content, i.e. Sperner and antichain bounds on the
%%       observation ladder. The antichain census needs a comparator.
%%   (ii) SCOPE: at this length it is a short-format result. Either expand it
%%       beyond the cube instance to the general case, or fold it into Paper 3,
%%       which is the other computational paper. Decide before submission; do not
%%       submit at this length with neither done.

% SIBLING PAPERS: Paper 1 (Obstruction Certificates under Incomplete Observation), Paper 2 (Exact Probabilistic Sufficiency), Paper 3 (Computational Certification of Obstruction), Paper 4 (Minimax Dual Certificates).
% Cross-cite, do not re-derive: the selector principle, the epistemic kernel
% and the observation structure are defined in Paper 1. Papers 2-5 cite it.

% SIBLING PAPERS: Paper 1 (Obstruction Certificates under Incomplete Observation), Paper 2 (Exact Probabilistic Sufficiency), Paper 3 (Computational Certification of Obstruction), Paper 4 (Minimax Dual Certificates).
% Cross-cite, do not re-derive: the selector principle, the epistemic kernel
% and the observation structure are defined in Paper 1. Papers 2-5 cite it.

% RESOURCED 2026-09-29: previous assembly (v10) was built from _v9 of this source,
%   4 versions behind. Re-assembled against paper2_exact_belief_computation_v13.tex. Content is a superset:
%   no labels dropped. See VERSION_CURRENCY.md.

\documentclass[10pt,twocolumn]{article}
\usepackage[a4paper,margin=15mm]{geometry}
\usepackage{amsmath,amssymb}
\usepackage{amsthm}
\theoremstyle{plain}
\newtheorem{theorem}{Theorem}
\newtheorem{proposition}[theorem]{Proposition}
\newtheorem{lemma}[theorem]{Lemma}
\newtheorem{corollary}[theorem]{Corollary}
\theoremstyle{definition}
\newtheorem{definition}[theorem]{Definition}
\theoremstyle{remark}
\newtheorem{remark}[theorem]{Remark}
\usepackage{booktabs,array,calc}
\usepackage[font=small,labelfont=bf]{caption}
\usepackage[colorlinks=true]{hyperref}
\providecommand{\tightlist}{\setlength{\itemsep}{0pt}\setlength{\parskip}{0pt}}
\emergencystretch=3em
\allowdisplaybreaks
\usepackage{graphicx}


\begin{document}
\title{Exactness is not the limitation: belief-state computation at scale and the audits that bind the obstruction calculus to its worked systems}
\begin{abstract}
\noindent\textbf{The general claim.} Exact arithmetic is widely treated as a constraint that
confines a theory to toy instances --- something one does on small examples before moving to the
floating-point treatment that is supposed to be the real contribution. This paper argues the
opposite on two independent fronts: \emph{exactness is the instrument, not the limitation.} It
scales further than the relaxing instinct expects, and it is precisely what makes the
counter-intuitive verdicts trustworthy. On the first front, the exact discipline survives the
curse of dimensionality: on a four-parameter cube of $256$ augmented cells with $17$ controls, the
stored object (an \emph{antichain}) and the argued instrument (a \emph{pairwise bound}) together
collapse $1{,}048{,}576$ subset evaluations to $496$ sets without loss, and at five parameters the
same two instruments carry $68{,}719{,}476{,}736$ evaluations onto $1{,}552$ sets. On the second
front, exact audits give every mechanism of the obstruction calculus a count, a witness and a
boundary --- and the verdicts they produce are not the ones a continuous or approximate treatment
would have predicted: a review-timing certificate that is \emph{non-strict}, holding everywhere
beyond $z_0 - 1$ and nowhere at the boundary; a monitoring-design target on which the adequate
partitions are enumerable but \emph{no coarsest adequate partition exists}; and two middle cells
that are simply incomparable. None of these is visible without exactness, and none of them is
checkable without it. \emph{A result that can be regenerated by script is a result a reader can
audit rather than a claim a reader must trust.}

\noindent\textbf{Why two parts.} Each part establishes one front. Part I takes the exact discipline
into the dimensionality regime where it is expected to fail and shows it does not, recovering
structure a pencil would never find. Part II applies one exact audit framework --- integer and
rational arithmetic throughout, every count regenerated by script --- to a family of worked systems
that instantiate the mechanisms of the calculus, and shows that exactness is what converts the
calculus from a set of sufficient conditions into a set of \emph{separating examples}. Part I
supplies the scale; Part II supplies the bite. Neither contains the other: a method that scales is
not thereby shown to produce surprising verdicts, and a collection of surprising exact examples is
not thereby shown to survive dimensionality.

\noindent\textbf{What is found.} Part I: a pairwise-Hamming lemma forces every blind survivable set
to consist of Hamming-adjacent cells, and no three cells of the four-cube are pairwise adjacent ---
so \emph{no blind policy of any period keeps three or more cells safe}. The bands are exact ($16$
maximal singletons at the floor's edge; exactly the $32$ Hamming-adjacent pairs at every level
above it). The observation ladder doubles: each additional exact probe of one parameter doubles the
conditional kernel mass, so three probes exhaust observation. A four-order antichain census
completes the classification; and at five parameters the pairwise bound \emph{stops being
sufficient}, which is the sharper news. Part II: the audit yields the non-strict timing identity on
a $93$-cell grid; $7$ of $15$ partitions adequate with two incomparable coarsest adequate
two-cell partitions; policy-class value dropping unequally across aggregation and protocol
restrictions with the middle cells incomparable, their kernels' meet institutional setwise, and
exactly one alternation-forced pair beyond both; unanimity costing a fixed identified set of
beliefs, sustained by exactly a $4 \times 4$ product of free votes; regime blindness losing exactly
the corner's pairs while no frozen regime loses any; and a continuous benchmark separating viable
from obstructed at $Y^{*} = 27/5$ with rational certificates on both sides.

\noindent\textbf{Structure.} Part I is the computation at scale and Part II the audit framework and
its worked systems. Each is self-contained and reproduced in full from its companion paper without
condensation, opening with that paper's own abstract. A cross-part conclusion follows.
\end{abstract}
\maketitle

\tableofcontents

\newpage

\part*{Part I --- Exact belief-state computation at scale}
\label{part:scale}
\addcontentsline{toc}{part}{Part I --- Computation at scale}

\maketitle

\begin{abstract}
This paper pushes the exact discipline of belief-state safety values into the curse-of-dimensionality regime: a four-parameter cube with
\(2^{4} = 16\) hidden-parameter cells on \(16\) stock levels ---
\(256\) augmented cells, \(17\) controls, every computation in exact
integer and rational arithmetic. The structure is classified
completely. A pairwise-Hamming lemma (every pair of distinct actions
sums, over any two parameter cells, to at most twice their agreement
count) forces every blind survivable set to consist of
Hamming-adjacent cells, and no three cells of the four-cube are
pairwise adjacent: no blind policy of any period keeps three or more
cells safe. The bands are exact: \(16\) maximal singletons at the
floor's edge, exactly the \(32\) Hamming-adjacent pairs at every level
above it (the matched--singly-mismatched alternation rises at net
\(\tfrac15\) per two steps with a \(\tfrac1{10}\) dip). The
observation ladder doubles: each additional exact probe of one
parameter doubles the conditional kernel mass ---
\(\tfrac1{16} \to \tfrac18 \to \tfrac14 \to \tfrac12 \to 1\) at the
edge, \(\tfrac18 \to \tfrac14 \to \tfrac12 \to 1 \to 1\) above it ---
so three probes exhaust observation above the edge. Exact point-based
evaluation is verified against definition-level enumeration at seven
rational beliefs, three levels, four horizons; the four-step alpha-set
deduplicates \(83{,}521\) action sequences to at most \(545\) distinct
outcome vectors. The antichain census is the headline: \(1{,}048{,}576\) raw subset evaluations collapse to \(496\) stored maximal sets. The instance is then enlarged to five parameters (Section~\ref{scale-fivecube}), where the classification changes qualitatively and is given completely: the maximal blind-survivable sets are the \(32\) radius-one Hamming balls at every level, plus \(80\) four-cell sets appearing from \(z_{0} \ge 1.3\), so that \(16 \times 2^{32} = 68{,}719{,}476{,}736\) raw subsets collapse to \(1{,}552\) stored sets. The raw lattice grows \(65{,}536\times\) against a \(3.13\times\) growth in the stored object, improving compression by four orders of magnitude; and the pairwise-Hamming instrument, sharp at four parameters, is shown necessary but no longer sufficient at five. A third instance of the theory's additive deadline law closes the paper (threshold \(1 + T/2\), zero probe excursion).
\end{abstract}

\textbf{Keywords:} exact computation; partially observable control;
zero-suppressed sets; survivable sets; Hamming structure; observation
design; antichain compression

\subsection{Introduction}\label{scale-introduction}

Two working disciplines govern the exact treatment of belief-state safety values: evaluate exactly where you probe, and store the antichain of maximal survivable sets rather than the lattice (Abaee, 2026, An obstruction calculus). This paper asks whether the exact discipline survives the curse of dimensionality, and what new structure appears there. The answer is affirmative, and the structure changes: the four-parameter cube is not a larger version of the two-parameter instance but a different regime. The cube is the paper's running instance; every structural claim below is instantiated on it.
Its survivable sets are governed by Hamming adjacency rather than by
coordinate matching; its observation ladder is geometric rather than
one-shot; and its raw subset lattice (\(2^{16}\) per level) is four
orders of magnitude larger than its stored antichain. All claims below
are regenerated by the paper's verification script in exact integer
and rational arithmetic, standard library only. The contributions are:
the pairwise-Hamming classification of the cube's survivable-set
structure (Section \ref{scale-hamming}); the exact observation-ladder
doubling and the point-based evaluation at four parameters (Sections
\ref{scale-ladder}--\ref{scale-census}); and the antichain census with the
audited-structure classification and the deadline instance (Sections
\ref{scale-glance}--\ref{scale-deadline}).

\subsection{The cube instance}\label{scale-instance}

The cube's local vocabulary (\(z\), \(z_{0}\), \(T\), the cell
structure) is disjoint from the companion papers' colliding letters;
the horizon \(T\) carries the same meaning as the certification
companion's \(J_{\mathcal{I}}(T)\). The state is \((z, \theta)\) with stock \(z\) on the grid
\(z_{0} \in \{1.0, 1.1, \dots, 2.5\}\) and hidden parameter
\(\theta = (\theta_{1}, \theta_{2}, \theta_{3}, \theta_{4}) \in
\{+1, -1\}^{4}\): \(16\) parameter cells on \(16\) levels, \(256\)
augmented cells. Controls are \(u \in \{+1,-1\}^{4} \cup \{(0,0,0,0)\}\)
(\(17\) actions), and the one-step drift is
\[
\dot z \;=\; d(u, \theta) \;=\; -\tfrac12
\;+\; \tfrac15\,\textstyle\sum_{i=1}^{4} u_{i}\theta_{i}:
\]
the matched action \(u = \theta\) rises at \(+\tfrac3{10}\); a single
mismatch gives \(-\tfrac1{10}\), two give \(-\tfrac12\), three give
\(-\tfrac9{10}\), the antimatched action \(-\tfrac{13}{10}\), and the
hold action \(-\tfrac12\) for every cell. The floor is \(z \ge 1\).
The sum over all sixteen cells of any action's drifts is \(-8\) per
step: the cube as a whole always declines, and the blind controller
declared here sees no observation of \(\theta\).

\subsection{The pairwise-Hamming classification}\label{scale-hamming}

\begin{lemma}[pair-sum bound and its survival cost]\label{scale-lem:pairsum}
Fix the dimension \(m \ge 1\), cells \(\theta, \theta' \in
\{+1,-1\}^{m}\) at Hamming distance \(h\), the action set
\(U = \{+1,-1\}^{m} \cup \{(0,\dots,0)\}\), the drift
\(d(u, \theta) = -\tfrac12 + \tfrac15\langle u, \theta\rangle\),
the floor \(z \ge 1\), and the blind class: one action sequence for
every parameter branch, of any period and possibly aperiodic. Then:
(i) for every \(u \in \{+1,-1\}^{m}\),
\(\langle u, \theta\rangle + \langle u, \theta'\rangle
\le 2\,(m - h)\), and the hold action gives \(0\) for every pair;
(ii) if a blind policy keeps both cells above the floor indefinitely,
then the pair's long-run average inner product satisfies
\(\liminf_{T \to \infty} \frac1T \sum_{t < T}
\bigl(\langle u_{t}, \theta\rangle + \langle u_{t},
\theta'\rangle\bigr) \ge 5\).
\end{lemma}

\begin{proof}
(i) The coordinates split into \(m - h\) agreements and \(h\)
disagreements; on an agreement \(u_{i}\theta_{i} = u_{i}\theta'_{i}
\le 1\) and on a disagreement the two terms cancel, so the sum is
\(2 \sum_{\text{agree}} u_{i}\theta_{i} \le 2(m-h)\); the hold
action gives \(0\) termwise. (ii) Survival of cell \(\theta\) means
\(z_{0} + \sum_{t < T} d(u_{t}, \theta) \ge 1\) for every \(T\);
dividing by \(T\) and passing to the limit, the Ces\`aro average of
the drifts has \(\liminf \ge 0\), and since
\(d(u_{t}, \theta) = -\tfrac12 + \tfrac15\langle u_{t},
\theta\rangle\), the \(\liminf\) of the average of
\(\langle u_{t}, \theta\rangle\) is at least
\(\tfrac52\); the same holds for \(\theta'\), and the two
\(\liminf\)s add. For a periodic policy the average is the exact
cycle average and no \(\liminf\) is needed.
\end{proof}

\begin{corollary}[the cube's pairs are Hamming-adjacent]
\label{scale-cor:hamming}
At \(m = 4\): (ii) forces the pair average to reach \(5\) while (i)
caps it at \(2(4 - h)\), so \(5 \le 2(4 - h)\), i.e.\
\(h \le \tfrac32\): every blind survivable pair has \(h \le 1\),
and since survivability passes to subsets, every pair inside a blind
survivable set is Hamming-adjacent. In general the argument gives
\(h \le m - \tfrac52\) (Remark~\ref{scale-rem:scope}).
\end{corollary}

\begin{lemma}[no adjacency triangles]\label{scale-lem:triangle}
No three distinct cells of \(\{+1,-1\}^{m}\) (any \(m\)) are
pairwise at Hamming distance \(1\). Hence every blind survivable set
at \(m = 4\) has at most two members, under every blind policy ---
any period, and no period at all. The bound (i) of
Lemma~\ref{scale-lem:pairsum} is verified for all \(120\) cell pairs and
all \(17\) actions at \(m = 4\) (and for all \(496\) pairs and
\(33\) actions at \(m = 5\)), and
Lemma~\ref{scale-lem:triangle} over all \(\binom{16}{3} = 560\) triples, by the verification script.
\end{lemma}

\begin{proof}
Flip coordinates so that the first cell is \(\mathbf{1}\) (flips
preserve all Hamming distances). A second cell at distance \(1\) is
\(\mathbf{1}\) with exactly one coordinate \(i\) flipped. A third
distinct cell at distance \(1\) from \(\mathbf{1}\) is
\(\mathbf{1}\) with exactly one coordinate \(j\) flipped: if
\(j = i\) it is the second cell, and if \(j \neq i\) the two differ
from each other in the two coordinates \(i, j\), so their Hamming
distance is \(2\). No triangle exists. (Equivalently: the hypercube
is bipartite by coordinate parity, so it has no odd cycle.)
\end{proof}

\begin{proposition}[bands and maximality]\label{scale-prop:bands}
A singleton \(\{\theta\}\) is survivable from every \(z_{0} \ge 1\)
(matched play rises at \(+\tfrac3{10}\)). A Hamming-adjacent pair
\(\{\theta, \theta'\}\) is survivable exactly from
\(z_{0} \ge 1.1\): the alternation \((\theta, \theta')\) gives each
member drifts \(+\tfrac3{10}, -\tfrac1{10}\) per two steps --- net
\(+\tfrac15\) with a within-cycle dip of \(\tfrac1{10}\) --- certified
at \(60\) steps, and from \(z_{0} = 1.0\) the first cell to act under
its mismatched action dips below the floor. The maximal survivable
sets are therefore exactly the \(16\) singletons at \(z_{0} = 1.0\)
and exactly the \(32\) Hamming-adjacent pairs at every
\(z_{0} \ge 1.1\), and nothing larger under any blind policy. An
exhaustive search over all \(5{,}219\) one-, two-, and three-periodic
policies with \(60\)-step certificates reproduces this classification
independently.
\end{proposition}

\begin{proof}
Singletons: matched play has drift \(+\tfrac3{10}\) from every
\(z_{0} \ge 1\). Pairs: let \(\{\theta, \theta'\}\) be
Hamming-adjacent and alternate \((\theta, \theta')\); each member
takes the matched drift \(+\tfrac3{10}\) and the mismatched drift
\(-\tfrac1{10}\) once per cycle, a net of \(+\tfrac15\) per two
steps with the cycle minimum \(z_{0} - \tfrac1{10}\) after the
mismatched step --- survivable exactly from \(z_{0} \ge 1.1\)
(the \(60\)-step certificates confirm the cycle arithmetic).
No pair from \(z_{0} = 1.0\): for any first action \(u\), if
\(u = \theta\) then \(\langle u, \theta'\rangle = 2\) and cell
\(\theta'\) falls to \(z_{0} - \tfrac1{10} < 1\) in one step, while
any other \(u\) drives at least one member below the floor in its
first step (the hold gives \(-\tfrac12\) to both); no repair exists
because the floor binds at every step. Nothing larger: any set of
three or more cells contains two at Hamming distance \(\ge 2\)
(Lemma~\ref{scale-lem:triangle}), and that pair violates
Corollary~\ref{scale-cor:hamming}, so no blind policy keeps it above the
floor --- and survivability passes to subsets, so the whole set fails
with it. Maximality is immediate: singletons extend to pairs from
\(1.1\), pairs extend to nothing.
\end{proof}

\begin{remark}[scope in dimension: the ball construction]
\label{scale-rem:scope}
The pairs-only census is the four-dimensional cube's, and the boundary
is sharp. In general dimension the pair argument of
Corollary~\ref{scale-cor:hamming} gives only \(h \le m - \tfrac52\); and
the constant matched action \(u = \mathbf{1}\) gives a cell at
Hamming distance \(k\) the constant drift
\(-\tfrac12 + \tfrac{m - 2k}{5}\), nonnegative exactly when
\(k \le \tfrac{2m-5}{4}\). Hence the Hamming ball of radius
\(r^{*}(m) = \lfloor (2m - 5)/4 \rfloor\) around any cell is
blind-survivable from every \(z_{0} \ge 1\):
\(r^{*}(4) = 0\), \(r^{*}(5) = r^{*}(6) = 1\),
\(r^{*}(7) = r^{*}(8) = 2\). At \(m = 5\) the radius-one ball ---
\(1 + 5 = 6\) cells, drifts \(+\tfrac12\) and \(+\tfrac1{10}\)
under constant matched play --- survives from the floor's edge, where
the four-cube keeps only singletons: the classification changes
qualitatively at \(m = 5\), and the maximal sets for \(m \ge 5\) are classified in Section~\ref{scale-fivecube} at \(m = 5\), and left open at \(m \ge 6\). All eight radii and the sign pattern are
verified by the verification script, with \(60\)-step simulations of the
radius-one ball at \(m = 5\) and the \(29\)-cell radius-two ball at
\(m = 7\).
\end{remark}

\begin{remark}[the crude instrument is not sufficient]
\label{scale-rem:crude}
The averaging identity used on smaller instances --- a survivable set
\(S\) needs \(\sum_{i} |\sum_{\theta \in S} \theta_{i}| \ge
\tfrac52\,|S|\) --- is not sufficient here: the four-set
\(\{++++, +++-, ++-+, +-++\}\) attains the bound with equality
(\(4 + 2 + 2 + 2 = 10\)) yet is not survivable, since it pairs cells
at Hamming distance \(2\). The pairwise instrument of
Lemma~\ref{scale-lem:pairsum} is the sharp one.
\end{remark}

\subsection{The observation ladder doubles}\label{scale-ladder}

\begin{proposition}[geometric ladder]\label{scale-prop:ladder}
Prior \(b_{0}\) is uniform on the sixteen cells; a probe of one
parameter refines the support to a subcube. The maximal conditional
kernel mass doubles with each probe:
\begin{align*}
\tfrac1{16} \to \tfrac18 \to \tfrac14 \to \tfrac12 \to 1
&\quad (z_{0} = 1.0),\\
\tfrac18 \to \tfrac14 \to \tfrac12 \to 1 \to 1
&\quad (z_{0} \ge 1.1),
\end{align*}
for zero, one, two, three, and four probes. Above the band's edge
three probes exhaust observation --- the support is then a
Hamming-adjacent pair, jointly survivable, and the fourth probe buys
nothing --- while at the edge the fourth probe is worth
\(\tfrac12\), the single-cell mass the pair's dip excludes.
\end{proposition}

\begin{proof}
The companion theory's antichain formula makes the kernel value of a
prior the maximal survivable-set mass it charges, and survivability is
hereditary downward --- a policy keeping a set alive keeps every
subset alive --- so the value at the uniform prior on a support
\(S\) is \(\max_{A \subseteq S,\ A
\text{ survivable}} |A|/|S|\), attained inside the support. By
Proposition~\ref{scale-prop:bands} the survivable subsets of a subcube are
singletons and Hamming-adjacent pairs, so a \(2^{j}\)-cell subcube
has value \(2/2^{j}\) from \(z_{0} \ge 1.1\) and \(1/2^{j}\) at
the edge. The supports after \(0, 1, 2, 3, 4\) probes of one
parameter are subcubes with \(j = 4, 3, 2, 1, 0\), giving the two
ladders exactly; the fourth edge probe buys the last cell because at
the edge only singletons survive and the fully probed support is then
itself the surviving cell. The script recomputes all ten values by
exhaustive survivable-subset enumeration within each support.
\end{proof}

\subsection{Exact point-based evaluation at four parameters}
\label{scale-pbvi}

The evaluation is point-based in the POMDP literature's sense
(Lovejoy, 1991), with the distinction that the values here are exact at
the probed beliefs, not bounds.

\begin{proposition}[exact alpha-set values]\label{scale-prop:pbvi}
The alpha-vector representation of the underlying theory (its
class-restricted rational alpha-vectors) evaluated at probed beliefs
returns the exact value: for horizons \(k = 1, 2, 3, 4\), three stock
levels, and seven random rational beliefs, the maximum over the
deduplicated alpha-set equals the maximum over definition-level
trajectory enumeration --- all \(84\) pairings equal. The four-step
alpha-set deduplicates \(17^{4} = 83{,}521\) action sequences to at
most \(545\) distinct outcome vectors per level (and to a single
vector at the top level's one-step horizon, where every action is
safe --- the silent region of the cube). Point-based selection
chooses where to evaluate, never what the value is.
\end{proposition}

\begin{proof}
Both sides are maxima of the same linear functional over the same
finite set of action sequences: the alpha-set is the image of the
sequence set under the map that records, for each parameter cell,
whether the sequence survives at the level; the maximum of a function
over a finite set equals the maximum over its image, and
deduplication removes only repeated vectors, never the attaining one.
Exactness of each entry is the definition-level simulation; no
discretization enters.
\end{proof}

\subsection{The antichain census}\label{scale-census}

\begin{proposition}[the antichain, not the lattice]\label{scale-prop:census}
The raw subset lattice carries \(2^{16} = 65{,}536\) subsets per
level --- \(16 \times 65{,}536 = 1{,}048{,}576\) over the grid ---
while the stored object is the antichain of maximal survivable sets --- the zero-suppressed
discipline of representing combinatorial set families by their maximal
elements (Minato, 1993):
\(16\) sets at the edge and \(32\) per level above it, \(496\) stored
sets in total. The per-level compression ratio is \(4{,}096\times\) at
the edge and \(2{,}048\times\) above it; and by
Lemma~\ref{scale-lem:triangle} nothing is lost --- the antichain is the
whole classification, not a summarization of one.
\end{proposition}

\begin{proof}
The raw count is \(16 \times (2^{16} - 1) = 1{,}048{,}560\) nonempty
subsets; the headline uses the full lattice
\(16 \times 2^{16} = 1{,}048{,}576\), the larger and therefore
conservative figure. The stored count is
\(16 + 15 \times 32 = 496\) by Proposition~\ref{scale-prop:bands}. The
antichain is lossless by the theory's kernel formula: the
value of any prior is the maximal survivable-set mass, read off the
stored sets alone.
\end{proof}

\subsection{The audited structure at a glance}\label{scale-glance}

\begin{table}[t]
\centering\footnotesize
\caption{The cube's audited structure, one row per layer; every entry is a
verification-script output in exact rational arithmetic.}
\label{scale-tab:glance}
\begin{tabular}{@{}l p{4.9cm}@{}}
\toprule
Layer & Exact statement \\
\midrule
Per-step drift & matched $+\tfrac3{10}$; $k$ mismatches: $-\tfrac12-\tfrac{2k}{5}$; hold $-\tfrac12$ \\
Cube accounting & drifts sum to $-8$ per step over the 16 cells \\
Singletons & $16$ survivable (exactly the floor's edge) \\
Pairs & exactly the $32$ Hamming-adjacent pairs; no other pair \\
Larger sets & none (three-set and four-set candidates fail) \\
Raw vs stored & $2^{16} = 65{,}536$ raw subsets per level, $496$ stored maximal sets \\
Observation ladder & ten probe schedules, values doubled at each rung \\
Deadline law & $z_0 \ge 1 + T/2$, $T = 0, \dots, 4$ (all cells) \\
\bottomrule
\end{tabular}
\end{table}

\subsection{A deadline instance}\label{scale-deadline}

\begin{proposition}[the deadline law on the cube]\label{scale-prop:deadline}
Hold for \(T\) steps (drift \(-\tfrac12\)), let the parameter be
revealed exactly at \(T\), then play matched forever (drift
\(+\tfrac3{10}\) on every branch). The companion theory's additive
deadline law applies with blind-drift cost \(d_{0} = \tfrac12\) and
zero probe excursion: viability holds exactly when
\[
z_{0} \;\ge\; 1 \;+\; \tfrac{T}{2}, \qquad T = 0, \dots, 4,
\]
verified by direct simulation on the grid and, independently, by a
mechanized proof (Section~\ref{scale-methods}). The excursion vanishes
because the learned act rises on every branch immediately --- the
law's \(e_{p}\) term prices the probe step the free revelation does
not need.
\end{proposition}

\begin{proof}
Necessity: under the hold every cell drifts \(-\tfrac12\), so at
revelation every branch sits at \(z_{0} - T/2\), and survival to
revelation forces \(z_{0} - T/2 \ge 1\). Sufficiency: the hold keeps
every branch at \(z_{0} - t/2 \ge 1\) for \(t \le T\) exactly
under the same inequality, and after revelation the matched action
rises at \(+\tfrac3{10}\) on every branch, so no later constraint
binds. This is the theory's additive deadline law at
\(d_{0} = \tfrac12\), \(e_{p} = 0\); the direct branchwise argument
needs no scalar-additivity hypothesis, and the simulation certifies
the identity cell by cell on the grid. The mechanized layer proves
the same threshold as an \emph{equivalence} --- viability holds iff
\(z_{0} \ge 1 + T/2\) --- over an arbitrary ordered field and with
\(T\) symbolic, so neither the grid nor the restriction
\(T \le 4\) is load-bearing for the argument.
\end{proof}

\medskip
\noindent\textbf{Reciprocal check (computational-certification
companion).} The deadline law's radius structure is certified
externally by the certification companion (Abaee, 2026, Computational viability certification): on this
instance's audited domain, a radius-\(\delta\) belief cell inherits
the verdict exactly for \(\delta \le z_{0} - (1 + T/2)\), with
\(\delta = 0\) recovering the stored per-state bound and the
boundary cells zero-slack; the joint check is archived with the
companion's verification record.

\subsection{Beyond four parameters: the five-cube, completely classified}
\label{scale-fivecube}

Remark~\ref{scale-rem:scope} fixes the ball threshold and records that the classification changes
qualitatively at \(m = 5\). This section carries the enlargement through: it gives every
maximal blind-survivable set of the five-cube, at every level, together with the census and
deduplication figures. Every claim below is settled by the safety game of
Section~\ref{scale-methods} --- a greatest-fixed-point computation on the capped non-negative
orthant, in exact integer arithmetic --- and not by an averaging bound. That distinction is
load-bearing, and Remark~\ref{scale-rem:crude} is only the first place it bites.

\subsubsection{The five-cube instance}

Take \(m = 5\): \(2^{5} = 32\) hidden-parameter cells over the same sixteen stock levels, so
\(512\) augmented cells against the four-cube's \(256\), and \(2^{5} + 1 = 33\) actions
against \(17\). The drift law and the floor are unchanged, so from
\(\langle u, \theta\rangle = m - 2\,\mathrm{dist}(u,\theta)\) the drift in tenths is
\[
D(h) \;=\; -5 + 2m - 4h \;=\; 5 - 4h \qquad (m = 5),
\]
namely \(+\tfrac12\) at distance \(0\), \(+\tfrac1{10}\) at distance \(1\), then
\(-\tfrac3{10}, -\tfrac7{10}, -\tfrac{11}{10}, -\tfrac32\) at distances \(2,3,4,5\), with
\(-\tfrac12\) for the hold. Two things change at once: the matched cell rises faster
(\(+\tfrac12\) against \(+\tfrac3{10}\)), and --- the qualitative change --- a cell at
Hamming distance \(1\) now \emph{rises}, where at \(m = 4\) it fell. That is the statement
\(r^{*}(4) = 0\), \(r^{*}(5) = 1\) of Remark~\ref{scale-rem:scope} read at the drift table.

\subsubsection{The maximal sets}

\begin{proposition}[the five-cube's maximal blind-survivable sets]
\label{scale-prop:five}
For the five-cube instance above, the maximal blind-survivable sets are:
\begin{enumerate}
\item the \(32\) Hamming balls of radius \(1\), each of \(1 + 5 = 6\) cells, at \emph{every}
level \(z_{0} \ge 1.0\);
\item from \(z_{0} \in \{1.0, 1.1, 1.2\}\), \emph{nothing else} --- \(32\) maximal sets per
level;
\item from \(z_{0} \ge 1.3\), additionally one orbit of four-cell sets under the
\(3{,}840\)-element automorphism group (XOR by any cell, then any permutation of the five
coordinates), of size \(80\), with pairwise-distance multiset \(\{1,1,1,1,2,2\}\) ---
\(32 + 80 = 112\) maximal sets per level.
\end{enumerate}
No set of five or more cells other than the balls is maximal, and no set of three or fewer
is maximal.
\end{proposition}

\begin{proof}
\textbf{Candidates.} By Corollary~\ref{scale-cor:hamming} generalised, survival of a pair needs
\(h \le m - \tfrac52\), i.e. \(h \le 2\) at \(m = 5\), so every survivable set is a clique of
the graph on the \(32\) cells with edges at distance \(\le 2\), and every such set lies in a
maximal clique. That graph has \(192\) maximal cliques: \(32\) of size \(6\) and \(160\) of
size \(4\).

\textbf{Size \(6\).} The \(32\) size-\(6\) maximal cliques are exactly the radius-\(1\) balls.
Under the constant matched action \(u = v\) on the ball \(B(v,1)\), the centre receives
\(D(0) = +5\) and each of the five neighbours \(D(1) = +1\), all strictly positive, so the
slack vector grows from \((L,\dots,L)\) at every level, \(L = 0\) included. All \(32\) are
therefore survivable from the floor's edge.

\textbf{Size \(4\).} The \(160\) size-\(4\) maximal cliques fall into exactly two orbits under
the automorphism group, \(80\) each: the \emph{tetrahedral} orbit, all six pairwise
distances equal to \(2\), and the orbit with distance multiset \(\{1,1,1,1,2,2\}\). The
tetrahedral orbit is \emph{not} survivable at any level --- the safety game returns an empty
winning region at the diagonal state for every \(L \in \{0,\dots,15\}\). The second orbit is
survivable exactly from \(L \ge 3\), i.e. \(z_{0} \ge 1.3\). Both verdicts are stable under
increasing the state cap from \(16\) to \(40\).

\textbf{Smaller sets.} All \(912\) cliques of size \(\le 3\) are contained in some
radius-\(1\) ball, and survivability is hereditary downward, so none is maximal. A
size-\(5\) clique would lie in a maximal clique of size \(\ge 5\), hence in a ball, hence is
non-maximal.
\end{proof}

\begin{remark}[the pairwise instrument is necessary but no longer sufficient]
\label{scale-rem:notsuff}
At \(m = 4\) the pairwise-Hamming bound is sharp: it excludes every pair at distance
\(\ge 2\), and everything it permits is survivable. At \(m = 5\) it is \emph{not} sharp. The
tetrahedral orbit has all six pairwise distances equal to \(2\), which meets the bound
\(h \le 2\) with equality, and no blind policy of any period keeps it above the floor. So
the pair-sum condition of Lemma~\ref{scale-lem:pairsum} is necessary and not sufficient from
dimension five onward, and the classification above cannot be recovered from pairwise data
alone. Remark~\ref{scale-rem:crude} records the same failure for the cruder averaging instrument
at \(m = 4\); the point here is that the sharper instrument inherits the failure one
dimension later.
\end{remark}

\subsubsection{The census at five parameters}

The raw subset lattice at \(m = 5\) carries \(2^{32} = 4{,}294{,}967{,}296\) subsets per
level --- \(16 \times 2^{32} = 68{,}719{,}476{,}736\) over the grid --- against the
four-cube's \(1{,}048{,}576\). By Proposition~\ref{scale-prop:five} the stored antichain is
\[
3 \times 32 \;+\; 13 \times 112 \;=\; 96 + 1{,}456 \;=\; 1{,}552
\]
maximal sets. Per level the compression is \(2^{32}/32 = 134{,}217{,}728\times\) at
\(z_{0} \le 1.2\) and \(2^{32}/112 \approx 3.83 \times 10^{7}\) above it; overall
\[
\frac{16 \times 2^{32}}{1{,}552} \;=\; \frac{2^{32}}{97} \;\approx\; 4.43 \times 10^{7},
\qquad\text{against}\qquad
\frac{1{,}048{,}576}{496} \;\approx\; 2{,}114 \quad (m = 4).
\]
The scaling is the point of the exercise. Passing from four parameters to five multiplies
the raw lattice by \(2^{16} = 65{,}536\) and the stored antichain by \(1{,}552/496 \approx
3.13\): the object one must \emph{not} store grows exponentially, the object one must store
grows by a factor a little over three, and the compression ratio improves by roughly
\(2.1 \times 10^{4}\). Exactness at five parameters is therefore not a constant factor more
expensive than at four; it is exponentially further out of reach of naive enumeration and
essentially no further out of reach of the antichain.

\subsubsection{Point-based evaluation at five parameters}

The four-step \(\alpha\)-set at \(m = 5\) deduplicates \(33^{4} = 1{,}185{,}921\) action
sequences to at most \(11{,}145\) distinct outcome vectors per level, against
\(17^{4} = 83{,}521 \to 545\) at \(m = 4\): the raw enumeration grows \(14.2\times\) and the
deduplicated object \(20.4\times\), so deduplication buys proportionally \emph{more}, not
less, as dimension rises. By level the count is \(945\) at \(z_{0} \le 1.2\), rising to
\(11{,}145\) at \(z_{0} = 2.5\). The two endpoint behaviours of the four-cube persist: the
one-step top-level collapse to a single outcome vector (\(33\) sequences, one vector, every
action safe in the silent region), and the edge behaviour (\(33\) sequences, \(33\) distinct
vectors, no collapse at all).

\subsubsection{What is not claimed at six parameters and beyond}

The ball construction gives \(r^{*}(6) = 1\), so the \(64\) seven-cell radius-one balls of
the six-cube are blind-survivable from every level by the same constant-matched-play
argument. What is \emph{not} determined at \(m = 6\) is whether they are maximal: the
distance-\(\le 3\) graph there has \(10{,}752\) maximal cliques with maximum clique size
\(12\), against the ball's \(7\), and resolving their survivability needs a safety game in
\(12\) dimensions, which the present apparatus does not settle. The five-cube classification
above is complete; the six-cube is open, and the gap between \(7\) and \(12\) is where the
next qualitative change, if there is one, would have to appear.

\subsection{Related work, and what this paper is}
\label{scale-priorart}

This paper is a computational companion, and the combinatorics it leans on is classical.
Both facts are stated here rather than left for a reader to discover.

\subsubsection{The place of this paper in the series}

The title's ``Scale~II'' refers to a companion relation, not to a volume number. The
theory being computed here is the belief-state safety value \(V_k(b)\) --- the maximal
probability of remaining safe for \(k\) steps under the disturbance's worst case --- together
with its support identity and \(\alpha\)-vector structure; that theory, and the
Sperner-bounded antichain formula behind its witness census, is developed in the companion
paper (Abaee, 2026, Belief-state safety values for viability under incomplete observation;
in this set, Paper~2). The present paper takes that theory's objects into the
curse-of-dimensionality regime on one explicit instance and classifies the instance
completely in exact arithmetic. It proves no new result about the theory itself. A reader
wanting the theorems should read that paper; a reader wanting the instance classified, and
the cost of doing so exactly, should read this one.

\subsubsection{Antichains, Sperner theory and Dedekind's problem}

An antichain in the Boolean lattice is a family of pairwise incomparable subsets. Sperner's
theorem (Sperner, 1927) gives the maximum size of an antichain on an \(n\)-set as
\(\binom{n}{\lfloor n/2\rfloor}\); counting \emph{all} antichains is Dedekind's problem,
with the asymptotic count \(2^{\binom{n}{\lfloor n/2\rfloor}(1+o(1))}\) due to Kleitman
(1966) and exact values known only for small \(n\) (Stanley, 2013; Engel, 1997, for the
standard treatments). Hamming adjacency, and the coding-theoretic machinery built on it, is
likewise classical.

\textbf{This paper proves no new theorem about antichains.} The census of
Section~\ref{scale-census} is a computation \emph{on one instance}: \(1{,}048{,}576\) raw subset
evaluations collapsing to \(496\) stored maximal sets. That figure is an output, and its
value is that it was obtained exactly and in full. The classical results are the scaffold,
not the contribution.

\subsubsection{Point-based and exact computation for partially observed models}

Point-based value iteration and its successors (Pineau, Gordon and Thrun, 2003; Shani,
Pineau and Kaplow, 2013, for the survey) approximate belief-space value functions by
propagating a finite set of belief points, with error bounds in the usual formulations. The
distinction here is not the idea of a point set but the arithmetic: every quantity in this
paper is an exact integer or rational, so the catalogue is not an approximation with a
bound but a classification of the instance. That costs feasibility at scale --- which is why
the instance is a four-parameter cube and not a larger one, and why
Section~\ref{scale-scope} says so. On the qualitative side, the complexity and decidability
landscape for partially observed models under qualitative objectives is established
(Chatterjee, Doyen and Henzinger, 2009); the companion paper's position relative to that
literature is stated there, and this paper adds nothing to it.

\subsubsection{What is claimed}

Given the above, the contribution is exactly five things.

\begin{enumerate}
\item \textbf{A complete classification of the four-parameter cube instance} --- every
  maximal blind survivable set, every band, and the observation ladder, in exact integer
  and rational arithmetic, with no approximation and no sampling.
\item \textbf{The cost figures for doing so.} The census collapse
  (\(1{,}048{,}576 \to 496\)), the four-step \(\alpha\)-set deduplication
  (\(83{,}521 \to 545\)), and the doubling of the observation ladder. These are the
  empirical content of the paper: they say what exactness costs at this size.
\item \textbf{The pairwise-Hamming structural restrictions as they apply here} --- that
  every blind survivable set consists of Hamming-adjacent cells, and that no three cells of
  the four-cube are pairwise adjacent, so no blind policy of any period keeps three or more
  cells safe. The adjacency notion is classical; its consequence for this instance is the
  result.
\item \textbf{The complete classification of the five-parameter cube}
  (Section~\ref{scale-fivecube}) --- the enlargement that Section~\ref{scale-scope} previously left open. The
  maximal sets are the \(32\) radius-one balls at every level, plus one \(80\)-set orbit of
  four-cell sets from \(z_{0} \ge 1.3\); the tetrahedral orbit meets the pairwise-Hamming bound
  with equality and is nevertheless not survivable, so that instrument is necessary and not
  sufficient from dimension five (Remark~\ref{scale-rem:notsuff}); and the census collapses
  \(68{,}719{,}476{,}736\) raw subsets to \(1{,}552\) stored sets, a compression
  \(2.1 \times 10^{4}\) times better than at four parameters.
\item \textbf{Machine-checking of the layer}, described next.
\end{enumerate}

\subsection{Mechanized verification}
\label{scale-mech}

The results above are additionally machine-checked. Sixteen modules of the project's
Lean~4 formalization are devoted to this paper's objects, in namespace
\texttt{Formalizations.EBC}: \texttt{EBC\_Bands}, \texttt{EBC\_Classification},
\texttt{EBC\_Deadline}, \texttt{EBC\_Dynamics}, \texttt{EBC\_ExactBelief},
\texttt{EBC\_Hamming}, \texttt{EBC\_Ladder} and \texttt{EBC\_Pairs}, with successive
versions of several. Together they carry \(134\) theorem and lemma declarations. The
mapping onto the sections here is by name: \texttt{EBC\_Hamming} carries the
pairwise-Hamming classification of Section~\ref{scale-hamming}, \texttt{EBC\_Bands} the exact
bands, \texttt{EBC\_Ladder} the doubling of Section~\ref{scale-ladder},
\texttt{EBC\_ExactBelief} the point-based evaluation of Section~\ref{scale-pbvi}, and
\texttt{EBC\_Deadline} the deadline instance of Section~\ref{scale-deadline}.

\noindent\textbf{No admitted gaps.} The full project builds cleanly under Lean~4 (\(60\) of \(60\) build jobs). In Lean an unfinished proof surfaces as the axiom
\texttt{sorryAx} in the axiom footprint of every declaration depending on it; a
source-level scan with comments stripped finds no \texttt{sorry}, \texttt{admit} or
\texttt{axiom} declaration anywhere in the formalization layer, all sixteen EBC modules
included. Declarations using classical reasoning depend only on \texttt{Classical.choice}, \texttt{propext} and \texttt{Quot.sound}, the standard axioms of Lean's classical logic.
\noindent\textbf{Verification provenance.} The project is pinned to Lean~4 in
\texttt{lean-toolchain} (currently \texttt{v4.34.1}). The mechanized checks reported in this
section were run under \texttt{v4.14.0}. On 2026-09-30 the source-level invariants were
re-verified against the current tree: a comment-stripped scan of all \texttt{.lean} files in
the project finds no \texttt{sorry}, \texttt{admit}, \texttt{axiom} or \texttt{constant}
declaration anywhere, and the declaration counts and named declarations quoted below were
re-counted and confirmed present. \textbf{The build has not been re-run since the toolchain
pin moved}, so the build-job figures below date from the \texttt{v4.14.0} run and should be
re-confirmed before submission.


\noindent\textbf{Scope.} What is machine-checked is the discrete combinatorial core --- the
structure the enumeration walks. The enumeration itself, including the census of
Section~\ref{scale-census}, is a computation performed by the deposited scripts and is
\emph{not} formalized.

\subsection{Verification methods}\label{scale-methods}

Every number is regenerated by one deterministic script in exact integer and rational arithmetic, standard library only. The battery opens with the pair-sum bound over all \(120\) pairs and \(17\) actions; the
adjacency-triangle check over all \(560\) triples; the exhaustive
one-, two-, three-periodic classification with \(60\)-step
certificates; the ten ladder values by survivable-subset enumeration
within each support; the \(84\) point-based pairings against
definition-level enumeration; the census counts; the deadline
thresholds by direct simulation; the dimension-scope battery (the
pair-sum bound at \(m = 5\) over all \(496\) pairs and \(33\)
actions; the ball radii \(r^{*}(m)\) for \(m = 3, \dots, 8\) with
the drift sign pattern and \(60\)-step simulations of the radius-one
ball at \(m = 5\) and the \(29\)-cell radius-two ball at
\(m = 7\)); and every headline number asserted present in this
text.

\textbf{A second, independent apparatus.} The script above enumerates;
it does not prove. Alongside it, the results of
Sections~\ref{scale-hamming}--\ref{scale-deadline} are formalized in a mechanized
layer (Lean~4, 60 modules under \texttt{Formalizations/}) with no axioms, no admitted
gaps, and no reliance on a computer-algebra oracle: every statement is
derived from the ordered-field interface alone, so it holds in any
model of that interface. The two apparatus were built separately and
agree where they overlap --- the ladder's first two entries
(\(1/16\) at the edge, \(1/8\) above it), the
adjacency-triangle fact behind the \(560\)-triple check, and the
\(16\)-singleton / \(32\)-pair census. Agreement of an enumeration
with a proof is evidence about both; it is not a substitute for either.

What each result now rests on, stated plainly:

\begin{itemize}
\item \textbf{Proved.} The band structure at the floor and above the
edge (Proposition~\ref{scale-prop:bands}); the ladder's value form
(Section~\ref{scale-ladder}); the deadline threshold
(Section~\ref{scale-deadline}). These hold as stated over an arbitrary
ordered field, with the support size and the horizon symbolic.

\item \textbf{Enumerated, not proved.} The census, the \(84\)
point-based pairings, the \(83{,}521 \to 545\) deduplication, and the
dimension-scope battery at \(m \ne 4\). These are finite checks over
the cube; their content is the count, and a count is what the script
supplies.

\item \textbf{Simulated.} The \(60\)-step ball certificates at
\(m = 5\) and \(m = 7\), and the grid certification of the deadline
thresholds. For the deadline the simulation is now redundant --- the
threshold is proved --- and is retained as an independent cross-check.
\end{itemize}

\noindent Where this text previously attributed a result to simulation
or enumeration alone, the attribution above governs.

\subsection{Delimitations}\label{scale-scope}

The blind class is declared and load-bearing: the policy may not
depend on the belief, so a single action sequence governs every
parameter branch and the pair-sum argument applies; a
\(z\)-feedback policy could branch its actions across branches (the
branches' stock paths diverge) and is outside the classification, as in the class discipline of the companion theory
(Abaee, 2026, An obstruction calculus), where the blind and
observation-dependent recursions are shown to differ --- the
feedback family contains the blind one, strictly, under
hypotheses of the same shape as this instance's. That result is
proved \emph{in the companion's model}; the correspondence with
the cube here is one of form, not of type, so it motivates
restricting to the blind class rather than proving that
restriction necessary. It does, however, fix the meaning of the
open question recorded above: since the two disciplines differ
strictly somewhere, whether a non-blind discipline beats the
pairs \emph{on the cube} cannot be settled by transferring the
companion's strictness, and is left open. Drifts are
deterministic per branch; probes are exact and cost no stock beyond
the blind decline they interrupt. The pairs-only classification is scoped to
\(m = 4\): in dimension five the constant matched action already
keeps the six-cell radius-one ball alive from every level
(Remark~\ref{scale-rem:scope}), and the maximal sets for \(m \ge 5\) are classified in Section~\ref{scale-fivecube} at \(m = 5\), and left open at \(m \ge 6\). Whether some non-blind discipline
beats the pairs on the cube is not claimed either way. Scaling beyond thirty-two cells --- where even the antichain's input enumeration becomes the bottleneck --- is the next step. Section~\ref{scale-fivecube} shows that the ball construction alone does not settle it: at \(m = 6\) the balls have seven cells while the maximal cliques of the pairwise-necessary graph reach twelve, and closing that gap needs a safety game the present apparatus does not carry.

\subsection{Conclusion}\label{scale-conclusion}

The exact discipline survives the curse of dimensionality and the
cube repays it with structure a pencil would never find: Hamming
governance of survivable sets, a doubling observation ladder, a
silent region at the top of the stock grid, and a four-order
antichain census. The stored object is the antichain; the argued instrument is the pairwise bound; and the two together are what make \(1{,}048{,}576\) subset evaluations collapse to \(496\) sets without loss. At five parameters the same two instruments, with the safety game in place of an averaging argument, carry \(68{,}719{,}476{,}736\) evaluations onto \(1{,}552\) sets --- and they carry the sharper news that the pairwise bound has stopped being sufficient.

\part*{Part II --- Exact audits of worked systems}
\label{part:audits}
\addcontentsline{toc}{part}{Part II --- Exact audits of worked systems}

\maketitle

\begin{abstract}
Viability certificates are useful in applications only insofar as their behaviour on concrete systems is understood precisely. This paper develops one exact audit framework --- integer and rational arithmetic throughout, every
count regenerated by script --- and applies it to a family of worked
systems that instantiate the mechanisms of the obstruction calculus
: a shared finite two-coordinate system with observation
structures (aggregation, policy-class restrictions, decentralized
observation, regime uncertainty), a review-timing grid for a
hidden-regime resource model, and a continuous two-floor production
benchmark. The audits yield the following. The review-timing identity on
the 93-cell grid is non-strict: viability at the deadline holds if and
only if \(T_{\mathrm{obs}} \le \tau^{*}(z_0) = z_0 - 1\), boundary included; the strict form against the real exit bound fails exactly on the boundary cells (against the floored count \(\sigma^{*}\), on one cell per row from \(z_0 = 2\)). On the four-state
monitoring target, 7 of the 15 partitions are adequate with two incomparable coarsest adequate two-cell partitions (minimal in the refinement order --- hence no single coarsest adequate design exists), and the two maximal common-action sets overlap. Four policy classes have kernel
sizes 24, 26, 25, and 28 of 36 belief pairs --- the classes form a two-by-two product of aggregation level and protocol --- with incomparable middle
cells. Decentralized unanimity over a
\(256\)-law space (\(64\) behaviorally distinct) sustains 12 of 36
pairs; a whole-law authority restriction collapses the kernel to zero,
and either agency's low vote forced against the target kills it. Under
frozen regimes no pair is lost (the corner pins its state); mode-blind
operation loses exactly the pairs containing the corner state
\((1,1)\). The biased certainty-equivalence law has
drift at least \(21/100\) on \([1,2]\); the bias-free law has drift
identically zero. The benchmark's critical aggregate is \(Y^{*} = 27/5\) (feasible,
margin \(0\)); at the audited obstructed aggregate \(Y = 6\) a
rational dual certificate prices the shortfall at margin \(3/50\). Static duality
holds at value \(-1/10\) on the two-floor instance with an \(\ell_1\)-normalized
Farkas pair, and a Helly-tight triple exhibits the sparse \(m{+}1\)-constraint witness. The complete result tables are carried here: the master table
of all 36 belief pairs against the five kernels, the sixteen sustaining
law pairs of the decentralized protocol (an exact \(4 \times 4\)
product), the 93-cell timing grid, the fifteen-partition census with
each fibre's common safe-action set, and the benchmark and drift tables,
with a five-figure set regenerated from the same scripts. Design rules
follow: review deadlines are boundary-inclusive; merged fibres call for
splits, not threshold tuning; aggregated indices support certainly-safe
reporting only; biased observation maps call for structural correction.
\end{abstract}

\textbf{Keywords:} viability kernels; obstruction certificates; exact
arithmetic; monitoring design; \mbox{decentralized observation}; certainty
equivalence; max-plus algebra

\subsection{Introduction}\label{audit-introduction}

The obstruction calculus states sound sufficient conditions for
nonviability --- obstruction certificates --- mechanism by mechanism:
common-action failure (Sections \ref{audit-master}--\ref{audit-adequacy}),
review-timing failure (Section \ref{audit-timing}), finite-time exit under
regime alternation (Section \ref{audit-regimes}), and failure induced by
restricting the policy class or by a biased observation map
(Sections \ref{audit-classes}, \ref{audit-drift}) (Abaee, 2026, An obstruction calculus). Each mechanism is
accompanied in application by quantitative questions. How many belief
states does an observation structure destroy, and which ones? At which
review deadlines does the timing certificate fire, and does it fire at
the boundary? What exactly does a vote protocol cost, and which law in
the protocol's space sustains a given target? This paper answers such
questions exactly for a family of worked systems, under one audit frame
and one notation. The sharp thesis, made explicit by the master
theorem of Section~\ref{audit-mtheorem}: obstructions arise from five
mechanisms --- common-action failure, memory/register coupling,
finite-time boundary exit, decentralized information loss, and joint
convex infeasibility --- and each admits a certificate with an explicit
semantic scope, a monotonicity status, and a robustness margin.

Three properties define the frame. First, all computations are carried
out in integer and rational arithmetic: no solver, no tolerance, no
simulation. Second, every discrete census in the family is finite --- the continuous objects reduce to exact rational calculations on rationally specified domains --- and every count is regenerated by a deterministic standard-library script, so any two executions agree. Third, each audited quantity is stated together with a
witness or a certificate --- an action sequence, a partition, a law
pair, a dual measure --- so the counts carry their own evidence. This paper carries the audit's complete result tables and its figure set:
the per-pair membership table behind the kernel counts, the sustaining
law pairs behind the decentralized count, the full timing grid, the full
partition census, and the benchmark and drift tables, every entry of
which is regenerated from the same scripts that verify the text.

The systems are drawn from the applied model of Abaee (2026, Robust viability of the 2J3KL limit reference point), in which
a two-coordinate resource stock is managed by acting one of two
instruments under incomplete observation. The family comprises: the
shared audit system with its observation structures (Sections
\ref{audit-master}--\ref{audit-decentralized} and \ref{audit-adequacy}--\ref{audit-regimes});
the hidden-regime review grid (Section
\ref{audit-timing}); a continuous two-floor production benchmark (Section
\ref{audit-benchmark}); and a multiple-floor line extending the Farkas
certificate to jointly binding constraints (Section \ref{audit-floors}). The
shared audit system is the running instance on which the paper's counts
are developed; each further member of the family is introduced against
it. The contributions are: the master monotonicity theorem --- one
audited-tuple object whose kernel comparisons unify the policy,
observation, and memory enlargements (Section \ref{audit-mtheorem}); the
complete exact audit of the worked family --- the master table, the
decentralized count, the full timing grid, the partition census, and the
benchmark and drift tables, every entry regenerated by the verification
scripts (Sections \ref{audit-master}--\ref{audit-regimes}); and the design rules
the audits license --- boundary-inclusive review deadlines,
merged-fibre caution, the monitoring-adequacy criterion, and the
five-mechanism obstruction taxonomy with its certificate semantics
(Sections \ref{audit-timing}--\ref{audit-adequacy}).

\subsection{Related work}
\label{audit-priorart}

\paragraph{Viability, and why worked systems.} That the governing question is which states
admit controls keeping every constraint satisfied over time, rather than which trajectory
optimises an objective, is the founding move of viability theory (Aubin, 1991). The obstruction
calculus audited here is the observation-constrained development of that programme, and the
systems below instantiate its mechanisms. Certificates are useful in application only insofar
as their behaviour on concrete systems is understood exactly, which is the purpose of this
paper: every count is regenerated by script in integer and rational arithmetic.

\paragraph{Worked benchmarks in the verification community.} The practice of maintaining a
shared, precisely specified family of benchmark problems on which competing methods are
compared is well established in the formal-verification community, where an annual friendly
competition applies a changing set of tools to a fixed set of benchmark problems in continuous
and hybrid systems, with results reported per benchmark rather than as a single ranking
(Althoff et al., 2020; Geretti et al., 2020). This paper's worked-systems family plays the
analogous role for the obstruction calculus: a fixed set of instances whose exact answers are
known and recomputable, on which the behaviour of the machinery can be read off.

The difference is what is being certified. The ARCH benchmarks are reachability problems for
continuous and hybrid dynamics, solved by reachability tools; the instances here are
\emph{information-constrained} viability problems under partial observation, and the quantities
tabulated are counts of survivable belief pairs, kernel sizes under policy-class restrictions,
and the sizes of minimal witness sets. No reachability tool computes these.

\paragraph{The classical instruments the audits exercise.} Three classical results recur and
are named where they are used: the alternative for systems of linear inequalities underlying
the duality certificate (Farkas, 1902); Helly's theorem on intersecting convex sets, which
governs the tightness of the sparse witness exhibited below (Helly, 1923); and max-plus
algebraic methods for the discrete-event structure of the review-timing grid (Baccelli, Cohen,
Olsder and Quadrat, 1992). None of these is contributed here; what is contributed is the exact
audit of what they do on each instance.

\paragraph{Exact and interval computation.} The audits run in integer and rational arithmetic
throughout. This is a statement about the verification pipeline and not a claim of superiority
over validated or interval-arithmetic approaches: the choice buys exactness at the cost of
scope, and the instances below are deliberately small enough to pay that price.

\subsection{Systems, conventions, and methods}\label{audit-frame}

\emph{The shared audit system.} States are pairs
\(x = (z_1, z_2) \in \{0,1,2,3\}^{2}\); the safe set is
\(\mathcal{V} = \{z_1 \ge 1,\ z_2 \ge 1\}\) with the two coordinates as
floors. Instruments are \(u \in \{0, 1, 2\}\): \(u = j \in \{1,2\}\)
acts coordinate \(j\), which regenerates one unit up to the cap \(3\),
while the other coordinate loses one unit; \(u = 0\) executes no
instrument and both coordinates lose one. The transition is defined index-wise: the acted coordinate regenerates, \(F_j(z)_j = \min\{3,\ z_j + 1\}\), and the other declines, \(F_j(z)_{3-j} = \max\{0,\ z_{3-j} - 1\}\), for \(j \in \{1,2\}\); \(F_0(z)_i = \max\{0,\ z_i - 1\}\) for both coordinates \(i\) --- coordinates clip at zero, the cap is saturating --- and safety is evaluated after each transition. A coordinate-swapped misreading (the regenerated coordinate written first for both \(j\)) reproduces no audited count and is excluded by the verification script. Instrument \(0\) is never uniquely viable on the audited family: admitting \(u = 0\) as a third action leaves all four kernel counts unchanged (\(24, 26, 25, 28\); verified by the verification script, which also checks the structural premises --- \(F_0 \le F_1\) and \(F_0 \le F_2\) pointwise on the grid and the safe set an up-set, the one-step reason no fibre needs \(0\)). The corner \((1,1)\) is the
single state from which no instrument maintains both floors, so its
singleton is nonviable; the nine safe states generate
\(C(9,2) = 36\) two-element belief pairs, the audited population (28 of the pairs full-information-viable; column F of Table~\ref{audit-tab:master}), at review horizon \(H = 12\) under the audited two-horizon verdict (the winning recursion at \(H\) and at \(H - 2\), which agree setwise on the audited universe and on the whole \(2^9\)-belief universe (stabilized from horizon \(2\) there, and from horizon \(1\) for the two non-alternation classes) --- \(W_{12} = W_{10}\) for each policy class, verified by the verification script --- so the verdict certifies a stabilized recursion; Abaee, 2026, An obstruction calculus, Theorem 1).

\emph{Observation structures.} An observation map
\(\mathrm{info}: x \mapsto c\) partitions a belief into fibres; a policy
acts per fibre on the read cell. The audited policy classes combine two
aggregation levels (full information; aggregation
\(c(x) = z_1 + z_2\)) with two protocols (free per-step choice; the
no-repeat protocol that forbids repeating the previous action. The
protocol's register is formalized as follows: one previous-action register
serves each review, shared by all fibres of the belief and not one per
branch; the policy fixes one action per fibre, the register then holds
some fibre's executed action, and the audited rule is pessimistic --- the
belief survives only if every admissible thread (each fibre's action
entering the register) leaves a winning continuation. Register values are
the two instruments, the initial register is unconstrained, and the survival condition quantifies over every
admissible thread, so the semantics is thread-independent by
construction (Proposition~\ref{audit-prop:register} computes what
single-thread rules and per-branch registers give instead).

\emph{Exactness.} Counts are computed by backward recursion over fibre
partitions with exact state arithmetic; witnesses are extracted as step-one per-fibre action vectors along a winning branch; the certificate itself is the recursion invariant --- winning-set membership at every horizon level, quantified universally over adversarial successor states and register threads --- of which the archived step-one vector is the first level. The census of
Section \ref{audit-adequacy} enumerates all \(15\) partitions of a four-state
target; partitions are ordered by refinement --- \(P\) is coarser than
\(P'\) when every fibre of \(P'\) lies in a fibre of \(P\) --- and
``coarsest adequate'' is relative to this order. The continuous objects of Section \ref{audit-benchmark} are rationals.

\emph{Conventions.} On the hold class of the hidden-regime model, the
worst-case exit time is \(\tau^{*}(z_0) = z_0 - 1\) and the discrete
survival supremum (the largest admissible integer review count) is
\(\sigma^{*}(z_0) = \lfloor z_0 - 1 \rfloor = \lfloor \tau^{*}(z_0)
\rfloor\), with level sets \(\{z_0 \ge 1 + k\}\); for an integer deadline
the inclusive criterion \(T_{\mathrm{obs}} \le \tau^{*}(z_0)\) coincides
with \(T_{\mathrm{obs}} \le \sigma^{*}(z_0)\), while the strict form
against the floored count is stricter than against the real exit time off
integer \(\tau^{*}\). Two
certainty-equivalence laws are audited: an uncorrected law whose index
leads the state by one step, \(s \mapsto s + 0.1\), and a corrected law
with the lead removed; both are read through the index \(g(s) = s^{2}\).

\emph{Figures and tables.} Every figure and every tabulated entry in
this paper is produced by a generation script that first recomputes the
quantity from the primitives above, in exact rational arithmetic, and
asserts it before drawing or printing; the master verification script of
Section \ref{audit-methods} chains both scripts and re-derives the tabulated
entries independently.

\medskip
\noindent\textbf{Shared symbols across the companion papers.} The
companion papers share this vocabulary with paper-local meanings; the
letters below carry meanings here that differ there:
\begin{center}\small
\begin{tabular}{@{}p{0.40\linewidth}p{0.52\linewidth}@{}}
\toprule
\textbf{Here} & \textbf{In the companions} \\
\midrule
\(\lambda\), \(\lambda^{*}\): the Farkas and dual certificate weights &
the Farkas multiplier and observer decay rate (obstruction calculus);
certificate weight vectors (certification; minimax duals) \\
\(\tau^{*}(z_0)\): the worst-case exit time &
observation and revelation times \(\tau\) (obstruction calculus;
certification; minimax duals) \\
\(\sigma^{*}(z_0)\): the floored exit count &
the kernel time variable \(\sigma\) (certification; minimax duals);
a noise scale (forecast ladder) \\
\bottomrule
\end{tabular}
\end{center}

\subsection{A master monotonicity theorem}\label{audit-mtheorem}

The audited family is an instance of one abstract object, and its
kernel comparisons are instances of one theorem. An \emph{audited
tuple} is
\(\Theta = (X, \mathcal{V}, A, F, \mathrm{info}, \Pi, m)\): a finite
state set \(X\) with safe set \(\mathcal{V}\), action set \(A\),
transition \(F\), observation map, admissible policy class \(\Pi\)
(per-cell action rules), and register convention \(m\) (shared, per-branch,
or none); viability is the horizon-stabilized recursion of
Section~\ref{audit-frame}.

\begin{proposition}[master monotonicity]\label{audit-prop:master}
Fix the audited tuple. Kernel membership is monotone in each of the
three enlargements, setwise: (i) \emph{policies}: if
\(\Pi \subseteq \Pi'\) then
\(\mathcal{W}_{\Pi} \subseteq \mathcal{W}_{\Pi'}\) --- restriction is
antitone (Proposition~\ref{audit-prop:lattice}); (ii) \emph{observations}:
if \(\mathrm{info}'\) is finer than \(\mathrm{info}\) (every fibre of
\(\mathrm{info}\) is a union of fibres of \(\mathrm{info}'\)) then
\(\mathcal{W}_{\mathrm{info}} \subseteq
\mathcal{W}_{\mathrm{info}'}\); (iii) \emph{memory}: a per-branch
register dominates the shared register --- every shared-register policy
is a per-branch policy, so the per-branch kernel contains the shared
one. Finer observation maps and larger memories can only enlarge the
kernel; all three statements are proved by the same induction on the
horizon: each enlargement enlarges the admissible action sets at every
cell and horizon level, so a winning continuation transfers upward.
\end{proposition}

\begin{proof}
Induction on \(h\). For \(h = 0\) the kernels coincide
(\(\mathcal{V}\) is fixed). Step: a belief winning at \(h\) under the
restricted tuple uses, at each cell, an action admissible and winning
there; the same choice is admissible under the enlarged tuple (the
action set at each cell only grows: more policies choose more actions,
finer cells shrink the quantifier over merged branches, a branch-local
register relaxes the no-repeat constraint on the other branch), and the
continuation is winning at \(h-1\) by the inductive hypothesis. The
finite population makes the induction an exact certificate: checks
P16--P18 verify the containments setwise on all \(36\) pairs.
\end{proof}

The audited counts are the theorem's rows: \(24 \le 26\) and
\(25 \le 28\) (observation axis, free protocol), \(24 \le 25\) and
\(26 \le 28\) (protocol axis at fixed aggregation), and the
per-branch-register kernels \((26, 26, 27)\) dominate the shared ones
\((24, 26, 25)\) (memory axis; Proposition~\ref{audit-prop:register}).
\begin{proposition}[stationary policies]\label{audit-prop:stationary}
Enumerate all deterministic stationary policies --- one action per cell,
fixed in time: \(2^{9} = 512\) at full information (nine singleton
cells), \(2^{5} = 32\) at aggregation (sums two to six). At full
information the union of the stationary policies' safe-pair sets is
exactly the kernel: memoryless sufficiency holds on the instance, and
the \(28\)-pair kernel admits an independent constructive description.
At aggregation the union is \emph{empty} --- no stationary policy
sustains any of the \(36\) pairs: the merged index is read against
states that demand conflicting actions across visits, and every
aggregate-viable pair relies on actions scheduled against time. The
\(26\)-pair aggregate kernel is therefore entirely a time-varying-policy
phenomenon --- the register-coupling mechanism in another guise, with
the review clock in the register's place (checks P17--P18; the fixpoint
computation of each class's horizon family independently reproduces all
four kernels).
\end{proposition}

\begin{proof}
Both counts are exhaustive enumerations with exact forward simulation
under each fixed policy (check P17), cross-validated against the
backward recursion by the horizon fixpoint (check P18).
\end{proof}

Strictness is instance-specific: the strict differences are exactly the
register-blocked pairs \(\{(1,2),(3,1)\}\) and \(\{(1,3),(2,1)\}\) on
the protocol axis, the single pair \(\{(1,3),(3,1)\}\) on the
observation axis, with the alternation-forced pair \(\{(1,2),(2,1)\}\)
outside both middle kernels; the incomparability of the
middle cells is a shared-register phenomenon, and the decentralized
kernel sits outside the chain because unanimity restricts \emph{who}
acts, not what the acting policy may do --- an orthogonal axis the
theorem prices separately (Section~\ref{audit-decentralized}).

\subsection{The master table}\label{audit-master}

Table \ref{audit-tab:master} is the audit's master table: all 36 belief pairs
against the five kernels --- the four policy classes of Section
\ref{audit-classes} (institutional: one shared compliance register per
review; aggregate: one pooled index per branch; full-information codex:
pooled registers across branches; full-information: unrestricted state
observation) and the decentralized kernel of Section
\ref{audit-decentralized}. Rows are grouped by the pair's lexicographically
smaller state, written \(z_1 z_2\).

\begin{table*}[t]
\centering
\footnotesize
\caption{The master table: membership of all 36 belief pairs in the five
audited kernels. I = institutional (aggregation, no-repeat); A =
aggregate; C = full-information codex; F = full information; D =
decentralized unanimity. A dot marks a non-member; every entry is
regenerated by script.}
\label{audit-tab:master}
\begin{tabular}{@{}l cccc c@{\hspace{12pt}}l cccc c@{}}
\toprule
Pair & I & A & C & F & D & Pair & I & A & C & F & D\\
\midrule
$12|21$ & $\cdot$ & $\cdot$ & $\cdot$ & F & D &
$22|31$ & I & A & C & F & D\\
$12|31$ & $\cdot$ & A & $\cdot$ & F & $\cdot$ &
$22|23$ & I & A & C & F & $\cdot$\\
$12|23$ & I & A & C & F & D &
$22|33$ & I & A & C & F & D\\
$12|33$ & I & A & C & F & $\cdot$ &
$22|32$ & I & A & C & F & $\cdot$\\
$12|22$ & I & A & C & F & $\cdot$ &
$23|31$ & I & A & C & F & $\cdot$\\
$12|32$ & I & A & C & F & D &
$23|33$ & I & A & C & F & $\cdot$\\
$12|13$ & I & A & C & F & $\cdot$ &
$23|32$ & I & A & C & F & D\\
$13|21$ & $\cdot$ & A & $\cdot$ & F & $\cdot$ &
$31|33$ & I & A & C & F & D\\
$13|31$ & $\cdot$ & $\cdot$ & C & F & D &
$31|32$ & I & A & C & F & $\cdot$\\
$13|23$ & I & A & C & F & $\cdot$ &
$32|33$ & I & A & C & F & $\cdot$\\
$13|33$ & I & A & C & F & D &
$11|12$ & $\cdot$ & $\cdot$ & $\cdot$ & $\cdot$ & $\cdot$\\
$13|22$ & I & A & C & F & D &
$11|21$ & $\cdot$ & $\cdot$ & $\cdot$ & $\cdot$ & $\cdot$\\
$13|32$ & I & A & C & F & $\cdot$ &
$11|31$ & $\cdot$ & $\cdot$ & $\cdot$ & $\cdot$ & $\cdot$\\
$21|31$ & I & A & C & F & $\cdot$ &
$11|23$ & $\cdot$ & $\cdot$ & $\cdot$ & $\cdot$ & $\cdot$\\
$21|23$ & I & A & C & F & D &
$11|33$ & $\cdot$ & $\cdot$ & $\cdot$ & $\cdot$ & $\cdot$\\
$21|33$ & I & A & C & F & $\cdot$ &
$11|22$ & $\cdot$ & $\cdot$ & $\cdot$ & $\cdot$ & $\cdot$\\
$21|22$ & I & A & C & F & $\cdot$ &
$11|32$ & $\cdot$ & $\cdot$ & $\cdot$ & $\cdot$ & $\cdot$\\
$21|32$ & I & A & C & F & D &
$11|13$ & $\cdot$ & $\cdot$ & $\cdot$ & $\cdot$ & $\cdot$\\
\bottomrule
\end{tabular}
\end{table*}

The column counts are the audited kernel sizes
\(|\mathcal{W}_{\mathrm{inst}}| = 24\),
\(|\mathcal{W}_{\mathrm{agg}}| = 26\),
\(|\mathcal{W}_{\mathrm{full\text{-}codex}}| = 25\),
\(|\mathcal{W}_{\mathrm{full}}| = 28\), and
\(|\mathcal{W}_{\mathrm{dec}}| = 12\). Figure \ref{audit-fig:kernels} displays
the four policy-class kernels, whose middle cells are incomparable. Two set-level facts are recorded exactly. The meet of the
two middle kernels equals the institutional kernel setwise. At class
level the join is the full-information class; at kernel level the union
of the two middle kernels falls one pair short of the full kernel:
exactly the pair \(\{(1,2),(2,1)\}\) lies in neither middle kernel. The
pair is the alternation-forced pair: under full information each branch
survives by its own strict alternation --- \((1,2)\) under \(1, 2, 1,
\dots\) and \((2,1)\) under \(2, 1, 2, \dots\) --- which the
idealized per-state policy can carry on both branches; aggregation cannot
see either requirement (one fibre, no common action), and the no-repeat
protocol, whose single previous-action register must serve both branches,
forbids the second branch's required next action once the first branch's
choice occupies the register. The full kernel is exactly the
\(C(8,2) = 28\) pairs avoiding the corner \((1,1)\): the corner's
singleton nonviability accounts for every remaining pair-level loss at
full information, as the eight \(11|\cdot\) rows of Table
\ref{audit-tab:master} show.

\begin{figure}[t]
\centering
\includegraphics[width=0.94\linewidth]{figs_ws4/fig_kernels.pdf}
\caption{The four policy classes form a two-by-two product of aggregation
level and protocol; the audited kernel sizes are \(24, 26, 25, 28\) of
36. The meet of the middle kernels is the institutional kernel setwise;
the class join is the full-information class, whose kernel is the 28
corner-free pairs, one pair beyond the middle-kernel union.}
\label{audit-fig:kernels}
\end{figure}

\subsection{Policy classes}\label{audit-classes}

The four audited classes combine full information and aggregation with
free and no-repeat protocols. Their kernels among the 36 pairs have
sizes:
\begin{gather*}
|\mathcal{W}_{\mathrm{inst}}| = 24,\quad
|\mathcal{W}_{\mathrm{agg}}| = 26,\\
|\mathcal{W}_{\mathrm{full\text{-}codex}}| = 25,\quad
|\mathcal{W}_{\mathrm{full}}| = 28 .
\end{gather*}
The four classes form a two-by-two product of aggregation level and
protocol; their kernels have incomparable middle cells:
the aggregate kernel contains pairs outside the full-information codex
kernel (for instance \(\{(1,2),(3,1)\}\) and \(\{(1,3),(2,1)\}\)), and
the codex kernel contains pairs outside the aggregate (for instance
\(\{(1,3),(3,1)\}\)). Within each observation column the free protocol's
kernel contains the no-repeat protocol's (both inclusions strict); across
columns neither middle kernel contains the other. The precise
meet-and-join reading is the one recorded in Section \ref{audit-master}. Two
per-pair readings illustrate the columns of Table \ref{audit-tab:master}. The
pair \(\{(1,2),(3,1)\}\) is aggregate-viable but not codex-viable: the
free protocol admits it under both observation structures, while the
no-repeat protocol excludes it under both --- at the second review the
thread that records the \((3,1)\)-branch's instrument \(2\) forbids the
forced instrument \(2\) of the successor branch \((2,1)\), the same
shared-register collision that excludes the alternation-forced pair. The pair \(\{(1,3),(3,1)\}\) is
codex-viable but not aggregate-viable: the two branches share the fibre
\(c = 4\), and no single per-fibre instrument serves both, while
full information separates them. Three structural facts behind the table
admit short proofs.

\begin{proposition}[full information decomposes branchwise]\label{audit-prop:decomp}
At full information under the free protocol, a two-state belief is
viable if and only if each of its singletons is viable; the kernel is
exactly the \(C(8,2) = 28\) corner-free pairs.
\end{proposition}

\begin{proof}
Full-information fibres are singletons, so a belief's step-one action
sets are the product of its branches': a choice is admissible exactly
when each branch's part is, and the recursion decomposes --- a belief is
viable exactly when both branches are. For the singletons: the corner
\((1,1)\) is nonviable because each instrument drops a coordinate to
zero in one step (\(u = 1\) forces \(z_2 = 0\), \(u = 2\) forces
\(z_1 = 0\), \(u = 0\) forces both); every other singleton is viable
--- the memoryless policy that plays instrument \(1\) exactly at
\(z_1 = 1\) and instrument \(2\) otherwise is safe forever: at
\(z_1 = 1\) the exclusion of the corner gives \(z_2 \ge 2\), so
\(u = 1\) restores \(z_1 = 2\) with \(z_2 \ge 1\); at
\(z_1 \ge 2\), \(u = 2\) keeps \(z_1 \ge 1\) and regenerates
\(z_2\).
\end{proof}

\begin{proposition}[class lattice]\label{audit-prop:lattice}
Kernel membership is antitone in the policy class: restricting the class
can only lose viable beliefs. Hence
\(\mathcal{W}_{\mathrm{inst}} \subseteq \mathcal{W}_{\mathrm{agg}},
\mathcal{W}_{\mathrm{full\text{-}codex}} \subseteq
\mathcal{W}_{\mathrm{full}}\); on the audited 36-pair population the
two middle inclusions are strict, the meet
\(\mathcal{W}_{\mathrm{agg}} \cap
\mathcal{W}_{\mathrm{full\text{-}codex}}\) equals the institutional
kernel setwise, and exactly one pair --- the alternation-forced pair
\(\{(1,2),(2,1)\}\) --- lies in the full kernel outside the middle
union.
\end{proposition}

\begin{proof}
Antitonicity --- order-reversal between class inclusion and kernel inclusion --- is immediate: a policy admissible for a restricted class is
admissible for any enlargements, so viability transfers upward. The
setwise facts are finite statements about a 36-element population; they
are proved by exhaustive certification --- the master table's
\(36 \times 5\) membership marks are recomputed row-by-row by the paper's verification script (checks M2 and M5), which is a complete
proof for a finite family.
\end{proof}

\begin{proposition}[register-semantics invariance]\label{audit-prop:register}
Endow the no-repeat protocol with the pessimistic shared-register
semantics: the policy fixes one action per fibre, the register then holds
some fibre's executed action, and the belief survives only if every
admissible thread leaves a winning continuation. Because the survival condition
quantifies universally over threads, the
pessimistic kernels depend on no thread choice: the audited counts
\(24, 26, 25, 28\) are thread-independent by construction. Two further
readings are computed for the record. A single-thread rule (one fixed
fibre's action entering the register) is order-sensitive: with the
first, or the last, fibre in sorted order the counts are
\(25, 26, 26, 28\) --- neither fixed rule yields the audited
institutional and codex counts --- so the well-defined semantics is the
pessimistic one. The optimistic semantics (some thread) is strictly
larger (codex \(27\)) and is not the audited reading. If instead each
branch carries its own register --- one compliance record per unit,
threaded forward along the branch --- the protocol cost collapses: the
institutional kernel equals the aggregate kernel setwise (\(26 = 26\)),
and the codex kernel becomes \(27\), regaining the alternation-forced
pair and losing exactly \(\{(1,3),(3,1)\}\) relative to full
information, with the aggregate kernel contained in it. The
incomparability of the middle cells is therefore a shared-register
phenomenon; the shared register is the audited institutional object ---
one compliance register per review across the units managed under the
protocol.
\end{proposition}

\begin{proof}
Every clause is a finite statement, computed from the shared primitives
of Section \ref{audit-frame} and asserted by the verification script: the
pessimistic counts (M2b), the optimistic containment (M2c), both
single-thread orders, and the per-branch kernels with the
\(\{(1,3),(3,1)\}\) loss (checks P9 and P10).
\end{proof}

\subsection{Decentralized observation}\label{audit-decentralized}

Two agencies each observe one coordinate and vote; an instrument
executes only on agreement. Each agency's vote is a map from its
coordinate's four values to the two instruments, so the law space has
\(2^{4} \times 2^{4} = 256\) law pairs. Unanimity sustains 12 of the
36 pairs (column D of Table \ref{audit-tab:master}), and exactly 16
coordinated-viable pairs are lost. The loss instance
\(\{(1,2),(3,1)\}\) is representative: each branch is individually
sustainable, and no law pair sustains the belief. For the audited belief
\(\{(1,2),(2,1)\}\), exactly 16 of the 256 law pairs sustain it; the
count coincides numerically with the 16 lost pairs, and no identity is
implied. The authority restriction is exact:

\begin{proposition}[authority restriction]\label{audit-prop:authority}
Fixing agency 2's vote at \(z_2 = 1\) to instrument \(1\) sustains no
law pair for the belief \(\{(1,2),(2,1)\}\): branch \((2,1)\) exits
at its first step under every such law pair. Fixing agency 1's vote at
\(z_1 = 1\) to instrument \(1\) leaves exactly the \(16\) sustaining
law pairs of Table \ref{audit-tab:laws}, and either whole-law
restriction (one agency's vote fixed to a single instrument at every
level) collapses the decentralized kernel to zero.
\end{proposition}

\begin{proof}
On branch \((2,1)\) at the first step: \(z_2 = 1\) must not
decline, so agreement on instrument \(2\) is required; agency 2's
fixed vote at \(z_2 = 1\) is \(1\), so agreement can hold only on
instrument \(1\), which drives \(z_2\) to zero, and disagreement
executes \(u = 0\), which does the same --- the branch exits at step
one under every admissible law pair. With agency 1's vote at \(z_1 = 1\) fixed to \(1\), the
sustaining pairs are exactly those with agency 1 forced to \(1\) at
\(z_1 = 1\) and \(2\) at \(z_1 = 2\) and agency 2 forced to \(2\)
at \(z_2 = 1\) and \(1\) at \(z_2 = 2\) (each forcing is a
one-step safety requirement: the low coordinate must regenerate), and
the free votes at \(\{0, 3\}\) never bind --- the \(4 \times 4\)
product of Table \ref{audit-tab:laws}. A whole-law restriction
fixes one agency's vote everywhere; the other agency's coordinate then
cannot regenerate at the low states where the fixed vote disagrees, and
the kernel collapses (verified exhaustively over all \(256\) law pairs, the whole-law
kernels \(0\) in all four cases; check P12 of the verification script).
\end{proof}

\begin{table}[t]
\centering
\footnotesize
\caption{The 16 law pairs sustaining the audited belief
\(\{(1,2),(2,1)\}\) factor exactly: agency 1's vote is forced to
\(1\) at \(z_1 = 1\) and \(2\) at \(z_1 = 2\), free at
\(z_1 \in \{0, 3\}\); agency 2's vote is forced to \(2\) at
\(z_2 = 1\) and \(1\) at \(z_2 = 2\), free at
\(z_2 \in \{0, 3\}\). Every combination sustains --- a complete
\(4 \times 4\) product (votes listed as
\((l(z{=}0), l(z{=}1), l(z{=}2), l(z{=}3))\));
agency 1's votes are the row labels, agency 2's the column labels.}
\label{audit-tab:laws}
\begin{tabular}{@{}l cccc@{}}
\toprule
 & \multicolumn{4}{c}{agency 2}\\
\cmidrule(l){2-5}
agency 1 & 1211 & 1212 & 2211 & 2212\\
\midrule
1121 & \checkmark & \checkmark & \checkmark & \checkmark\\
1122 & \checkmark & \checkmark & \checkmark & \checkmark\\
2121 & \checkmark & \checkmark & \checkmark & \checkmark\\
2122 & \checkmark & \checkmark & \checkmark & \checkmark\\
\bottomrule
\end{tabular}
\end{table}

The law space itself is counted before quotienting: a \(z = 0\) vote is
consultable only on a trajectory that has already left the safe set, so
the audited kernel is invariant under the \(z{=}0\)-vote quotient ---
the effective law space is \(2^{3} \times 2^{3} = 64\) behaviorally
distinct law pairs, and the \(16\) sustaining law pairs of Table
\ref{audit-tab:laws} comprise \(4\) distinct behavioral laws times the \(4\)
never-consulted \(z{=}0\) votes. The sustaining set has an
exact product structure (Table
\ref{audit-tab:laws}): on the trajectory alternating between the branches,
agency 1 is read only at \(z_1 \in \{1, 2\}\) and must vote \(1\) then
\(2\); agency 2 is read only at \(z_2 \in \{2, 1\}\) and must vote
\(1\) then \(2\); the votes at the unreachable coordinates
\(z \in \{0, 3\}\) are free, giving \(2^{2} \times 2^{2} = 16\). The
restricted kernels are exact. Fixing either agency's low vote to the
regenerating instrument (agency 1 to \(1\) at \(z_1 = 1\), agency 2 to
\(2\) at \(z_2 = 1\)) leaves the decentralized kernel untouched at
\(12\) --- the constraint does not bind. Fixing either agency's low
vote against it (agency 2 to \(1\) at \(z_2 = 1\), agency 1 to \(2\)
at \(z_1 = 1\)) shrinks the kernel to four beliefs ---
\(\{(2,3),(3,2)\}\), \(\{(2,2),(3,3)\}\), \(\{(1,3),(3,3)\}\),
\(\{(1,3),(2,2)\}\), up to the symmetry below --- and kills the
audited target; a whole-law restriction collapses the kernel to \(0\)
in all four cases. The model's swap symmetry (coordinates and instruments
together) interchanges the agencies and maps the audited target to
itself, so no agency is privileged: what carries the target is the
direction of each forced vote --- each branch's low coordinate must be
allowed to regenerate --- not the agency holding it. Two modelling choices are declared. The deadlock transition
\(u = 0\) executes no instrument and both coordinates lose one; and
each agency votes on its own coordinate's current value: the audited
protocol is memoryless and decentralized --- no agency observes the
belief, the other agency's observation, or anything like common
knowledge. The count has a theorem behind it:

\begin{proposition}[the decentralized kernel factors statewise]
\label{audit-prop:req}
Unanimity executes statewise: the executed instrument at a state depends
only on the two votes at that state's coordinates, and safety is
evaluated per state. Hence a belief is sustained if and only if its
branches share a common sustaining law pair,
\(\bigcap_{x \in B} S(x) \neq \emptyset\), where \(S(x)\) is the set
of law pairs sustaining the singleton \(\{x\}\); on the audited
population the factored criterion reproduces the joint kernel exactly
(the \(12\) sustained beliefs of the master table), and the \(16\)
sustaining law pairs of the audited belief
(Table~\ref{audit-tab:laws}) are invariant under the model's coordinate-swap
symmetry. The obstruction is legible branchwise: a belief is lost
exactly when its branches' sustaining sets do not intersect --- as for
the representative loss \(\{(1,2),(3,1)\}\), each branch individually
sustainable, no law pair sustaining both.
\end{proposition}

\begin{proof}
The belief trajectory under a law pair is the union of the branch
trajectories (statewise execution, statewise safety), so it stays in
\(\mathcal{V}\) if and only if every branch trajectory does --- which is
the displayed intersection. Both computations (the factored singleton
intersection and the joint \(256\)-pair enumeration) are run on all
\(36\) beliefs and agree setwise (check P20). The symmetry is the swap
map of this section's closing remark applied to the product structure.
\end{proof}

\subsection{Review timing on the hidden-regime grid}\label{audit-timing}

The hidden-regime resource model: a single stock held at level
\(z_0 \ge 1\) above the floor; each review period the unobserved
adversarial regime realizes one unit of decline and there is no
regeneration (the hold class), so the stock tracks \(z(t) = z_0 - t\)
and remains safe exactly while \(z(t) \ge 1\). The model is audited on
the grid \(z_0 \in \{1.0, 1.1, \dots, 4.0\}\) and
\(T_{\mathrm{obs}} \in \{1, 2, 3\}\), 93 cells. Viability at the
deadline holds if and only if
\(T_{\mathrm{obs}} \le \tau^{*}(z_0) = z_0 - 1\), with zero
mismatches across the grid. The identity is boundary-inclusive: the
viable cells number \(33\), and the strict form \(T_{\mathrm{obs}} <
\tau^{*}(z_0)\) fails exactly on the boundary cells \((2,1)\),
\((3,2)\), and \((4,3)\) --- those with \(z_0 = 1 +
T_{\mathrm{obs}}\); against the floored count \(\sigma^{*}\) the
strict form fails on one cell per row \(z_0 \ge 2\) --- \(21\) cells,
the boundary cells included --- which is why the
inclusive predicate is the audited identity.
Table \ref{audit-tab:timing} carries the full grid and Figure
\ref{audit-fig:timing} displays it. The practical form of the rule: a review
must land no later than the worst-case exit time, and a deadline exactly at the worst-case exit is admissible under the audited endpoint convention (closed safe set; safety evaluated at the deadline review); no open-set or alternative-exit-convention claim is made. The audited identity is in fact the
restriction of a general law:

\begin{proposition}[review timing, general form]\label{audit-prop:timing}
On the hold class of the hidden-regime model, for every initial level
\(z_0 \ge 1\) and every integer deadline \(T \ge 0\), the cell
\((z_0, T)\) is deadline-viable if and only if \(T \le z_0 - 1\).
Accordingly \(\sigma^{*}(z_0) = \lfloor z_0 - 1 \rfloor\) with
superlevel sets \(\{z_0 \ge 1 + k\}\), \(k \in
\mathbb{Z}_{\ge 0}\), and the audited 93-cell grid is the restriction
to \(z_0 \in \{1.0, \dots, 4.0\}\), \(T \in \{1, 2, 3\}\).
\end{proposition}

\begin{proof}
Under every admissible hold-class control the adversarial regime realizes
the declining branch at each step, so the worst stock declines at unit
rate and remains safe exactly through time \(z_0 - 1\) --- the last
safe review, the first violation occurring at \(z_0\):
deadline-viability forces \(T \le z_0 - 1\). Conversely
the hold policy tracks the decline path \(z(t) = z_0 - t\) branchwise,
which stays admissible exactly on \([0,\, z_0 - 1]\) and terminates at
the floor, which is safe because the safe set is closed; hence
\(T \le z_0 - 1\) is viable, boundary included. The set identity
\(\{T : T \le z_0 - 1\} = \{T : T \le \sigma^{*}(z_0)\}\)
(integer \(T\), all \(z_0\)) --- equivalently \(\{T < z_0 - 1\} =
\{T \le \sigma^{*}(z_0)\}\) off the integers, with the boundary
failure at integer \(z_0 - 1\) --- gives the displayed supremum.
\end{proof}

\begin{proposition}[hitting times on regime graphs]\label{audit-prop:regime}
Let the adversary choose a regime path \(r_1 r_2 \dots\) through a
finite regime graph with per-step damages \(d_r \ge 0\), subject to
the graph's adjacency constraint. The worst cumulative damage by time
\(t\) is the max-plus recursion
\(D_t = \max_{\text{paths}} \sum_{s \le t} d_{r_s}\) --- the standard
timed-event-system algebra (Baccelli et al., 1992) --- and the cell
\((z_0, T)\) is deadline-viable if and only if \(D_T \le z_0 - 1\)
(the audited endpoint convention). Two exact instances: with free
switching the adversary plays the worst regime every step, \(D_t = t
\cdot \max_r d_r\), recovering Proposition~\ref{audit-prop:timing} on the
audited grid; with two regimes of damages \(1\) and \(1/2\) forced to
alternate, the worst damage is \(D_{2k} = 3k/2\) and \(D_{2k+1} = 3k/2
+ 1\), and the re-audited \(93\)-cell grid has exactly \(43\) viable
cells (\(21, 16, 6\) at \(T = 1, 2, 3\)), zero mismatches between the
closed form and path enumeration. The max-plus dynamic program equals
path enumeration exactly on a three-regime graph with a forbidden
self-loop (all horizons to six).
\end{proposition}

\begin{proof}
Deadline-viability is safety of the worst adversarial path to the
deadline, which is the max-plus evaluation of the damage sum; free
switching decouples the steps, giving \(t \cdot \max_r d_r = t\) on
the audited damages and hence the audited identity. Under alternation
the two interleavings give \(D_{2k} = k(1 + \tfrac12) = 3k/2\) and
\(D_{2k+1} = (k+1) + \tfrac{k}{2}\); the grid recomputation enumerates
both interleavings per cell (check P21). The dynamic-programming (DP) evaluation over
(last regime, step) is exact rational arithmetic and is asserted equal
to explicit path enumeration.
\end{proof}

\begin{table}[t]
\centering
\footnotesize
\caption{The timing grid, printed in full (\(93\) cells): viability at the deadline against \(\sigma^{*}(z_0) = \lfloor
\tau^{*}(z_0) \rfloor\) (the deadline count; the identity columns use
the inclusive predicate, on which the floored and real-valued forms
coincide). The checkmark of column
\(T_{\mathrm{obs}}\) means the cell is viable; the identity
column ``id.'' is the predicate
\(T_{\mathrm{obs}} \le \sigma^{*}(z_0)\) and column ``cell'' the audited
viability; the two agree on all \(93\) cells (zero mismatches; \(33\)
viable cells). The circled cell of Figure \ref{audit-fig:timing} is
\((2.0, 1)\): viable, on the boundary; the strict form disagrees
exactly on the boundary cells \((2,1)\), \((3,2)\), \((4,3)\).}
\label{audit-tab:timing}
\begin{tabular}{@{}l c c cc c cc c@{}}
\toprule
 & & \multicolumn{2}{c}{$T_{\mathrm{obs}} = 1$} &
\multicolumn{2}{c}{$T_{\mathrm{obs}} = 2$} &
\multicolumn{2}{c}{$T_{\mathrm{obs}} = 3$}\\
\cmidrule(lr){3-4}\cmidrule(lr){5-6}\cmidrule(lr){7-8}
$z_0$ & $\sigma^{*}$ & id. & cell & id. & cell & id. & cell\\
\midrule
1.0 & 0 & $\cdot$ & $\cdot$ & $\cdot$ & $\cdot$ & $\cdot$ & $\cdot$\\
1.1 & 0 & $\cdot$ & $\cdot$ & $\cdot$ & $\cdot$ & $\cdot$ & $\cdot$\\
1.2 & 0 & $\cdot$ & $\cdot$ & $\cdot$ & $\cdot$ & $\cdot$ & $\cdot$\\
1.3 & 0 & $\cdot$ & $\cdot$ & $\cdot$ & $\cdot$ & $\cdot$ & $\cdot$\\
1.4 & 0 & $\cdot$ & $\cdot$ & $\cdot$ & $\cdot$ & $\cdot$ & $\cdot$\\
1.5 & 0 & $\cdot$ & $\cdot$ & $\cdot$ & $\cdot$ & $\cdot$ & $\cdot$\\
1.6 & 0 & $\cdot$ & $\cdot$ & $\cdot$ & $\cdot$ & $\cdot$ & $\cdot$\\
1.7 & 0 & $\cdot$ & $\cdot$ & $\cdot$ & $\cdot$ & $\cdot$ & $\cdot$\\
1.8 & 0 & $\cdot$ & $\cdot$ & $\cdot$ & $\cdot$ & $\cdot$ & $\cdot$\\
1.9 & 0 & $\cdot$ & $\cdot$ & $\cdot$ & $\cdot$ & $\cdot$ & $\cdot$\\
2.0 & 1 & \checkmark & \checkmark & $\cdot$ & $\cdot$ & $\cdot$ & $\cdot$\\
2.1 & 1 & \checkmark & \checkmark & $\cdot$ & $\cdot$ & $\cdot$ & $\cdot$\\
2.2 & 1 & \checkmark & \checkmark & $\cdot$ & $\cdot$ & $\cdot$ & $\cdot$\\
2.3 & 1 & \checkmark & \checkmark & $\cdot$ & $\cdot$ & $\cdot$ & $\cdot$\\
2.4 & 1 & \checkmark & \checkmark & $\cdot$ & $\cdot$ & $\cdot$ & $\cdot$\\
2.5 & 1 & \checkmark & \checkmark & $\cdot$ & $\cdot$ & $\cdot$ & $\cdot$\\
\midrule
2.6 & 1 & \checkmark & $\cdot$ & $\cdot$\\
2.7 & 1 & \checkmark & $\cdot$ & $\cdot$\\
2.8 & 1 & \checkmark & $\cdot$ & $\cdot$\\
2.9 & 1 & \checkmark & $\cdot$ & $\cdot$\\
3.0 & 2 & \checkmark & \checkmark & $\cdot$\\
3.1 & 2 & \checkmark & \checkmark & $\cdot$\\
3.2 & 2 & \checkmark & \checkmark & $\cdot$\\
3.3 & 2 & \checkmark & \checkmark & $\cdot$\\
3.4 & 2 & \checkmark & \checkmark & $\cdot$\\
3.5 & 2 & \checkmark & \checkmark & $\cdot$\\
3.6 & 2 & \checkmark & \checkmark & $\cdot$\\
3.7 & 2 & \checkmark & \checkmark & $\cdot$\\
3.8 & 2 & \checkmark & \checkmark & $\cdot$\\
3.9 & 2 & \checkmark & \checkmark & $\cdot$\\
4.0 & 3 & \checkmark & \checkmark & \checkmark\\
\bottomrule
\end{tabular}
\end{table}

\begin{figure}[t]
\centering
\includegraphics[width=0.96\linewidth]{figs_ws4/fig_timing.pdf}
\caption{The timing grid. Cells viable at the deadline (green) satisfy \(T_{\mathrm{obs}} \le \tau^{*}(z_0)\); the identity is
boundary-inclusive, and the strict form \(T_{\mathrm{obs}} <
\tau^{*}(z_0)\) fails at exactly one cell,
\((z_0, T_{\mathrm{obs}}) = (2.0, 1)\) (circled); against the floored count \(\sigma^{*}(z_0) = \lfloor z_0 - 1 \rfloor\) the strict form fails on all six viable cells, which is why the inclusive predicate is the audited identity.}
\label{audit-fig:timing}
\end{figure}

\subsection{Monitoring adequacy}\label{audit-adequacy}

Let \(R(x)\) denote the safe actions of state \(x\): those keeping
\(x\) in \(\mathcal{V}\) for one step. On the audited target,
\(R((1,2)) = \{1\}\), \(R((2,1)) = \{2\}\), \(R((2,2)) = \{1,2\}\), and
\(R((3,1)) = \{2\}\). A partition of a target belief is \emph{adequate}
when every fibre has a nonempty common safe-action set. On the
four-state target \(\{(1,2), (2,1), (2,2), (3,1)\}\), 7 of the 15
partitions are adequate, and two minimal adequate partitions with two
cells exist. Table \ref{audit-tab:census} carries the complete census with
each fibre's common safe-action set; Figure \ref{audit-fig:census} displays
it.

\begin{table}[t]
\centering
\footnotesize
\caption{The fifteen-partition census on the target
\(\{(1,2), (2,1), (2,2), (3,1)\}\) (states written \(z_1 z_2\); fibres
separated by ``;''). A partition is adequate when every fibre's common
safe-action set (brackets; \(\cdot\) = empty) is nonempty. Seven
partitions are adequate; the minimal adequate two-cell partitions are
rows 6 and 7.}
\label{audit-tab:census}
\begin{tabular}{@{}l l l@{}}
\toprule
Partition & Adequate & Common safe actions\\
\midrule
$12|21|22|31$ & no & $[\cdot]$\\
$12|21|31$ ; $22$ & no & $[\cdot\ 12]$\\
$12|22|31$ ; $21$ & no & $[\cdot\ 2]$\\
$12|31$ ; $21|22$ & no & $[\cdot\ 2]$\\
$12|31$ ; $22$ ; $21$ & no & $[\cdot\ 12\ 2]$\\
$21|22|31$ ; $12$ & \textbf{yes} & $[2\ 1]$\\
$21|31$ ; $12|22$ & \textbf{yes} & $[2\ 1]$\\
$21|31$ ; $22$ ; $12$ & yes & $[2\ 12\ 1]$\\
$22|31$ ; $12|21$ & no & $[2\ \cdot]$\\
$22|31$ ; $21$ ; $12$ & yes & $[2\ 2\ 1]$\\
$31$ ; $12|21|22$ & no & $[2\ \cdot]$\\
$31$ ; $12|22$ ; $21$ & yes & $[2\ 1\ 2]$\\
$31$ ; $21|22$ ; $12$ & yes & $[2\ 2\ 1]$\\
$31$ ; $22$ ; $12|21$ & no & $[2\ 12\ \cdot]$\\
$31$ ; $22$ ; $21$ ; $12$ & yes & $[2\ 12\ 2\ 1]$\\
\bottomrule
\end{tabular}
\end{table}

\begin{figure}[t]
\centering
\includegraphics[width=0.98\linewidth]{figs_ws4/fig_census.pdf}
\caption{The census. Each panel draws one partition: a bar over a group
of states is a fibre (green when its common safe-action set is
nonempty, red when empty), and the margin marks adequacy
(\(+\)/\(\times\)). The two adequate two-cell partitions are panels 6
and 7 of the middle row.}
\label{audit-fig:census}
\end{figure}

The maximal subsets with a nonempty common safe-action set are
\(\{(1,2), (2,2)\}\) and \(\{(2,1), (2,2), (3,1)\}\). They overlap in
\((2,2)\), so no partition into maximal common-action sets exists:
coarsening to maximal fibres is not an available design, and the
adequate partitions found by the census are the design options. The criterion is one-step
monitoring adequacy; dynamic viability is not claimed from it (the
kernels are computed by the full recursion). The census excludes
instrument \(0\) without loss: \(F_0 \le F_1\) and \(F_0 \le F_2\)
pointwise with the safe set an up-set (verified in the frame section)
give every one-step-safe fibre a safe instrument in \(\{1,2\}\). Two
consequences follow. Pairwise intersection of safe-action sets is necessary for fibre
adequacy; in the two-instrument audited universe it is also sufficient
(Helly number two: a pairwise-intersecting family of subsets of a
two-point action set has a common action), so every inadequate fibre in
the census fails with a literal empty intersection, while in action sets
of three or more elements pairwise intersection ceases to suffice
(\(\{1,2\}, \{2,3\}, \{1,3\}\)); and an
observation map merging states with disjoint safe-action sets obstructs
at that information state while every branch remains individually
viable --- the remedy is a fibre split, not a finer threshold on the
same index. The census makes both visible at the level of individual
fibres: in each inadequate partition of Table \ref{audit-tab:census} the
failing fibre is exactly one merging \((1,2)\) --- whose safe-action
set is the singleton \(\{1\}\) --- with a state whose safe actions
exclude \(1\), i.e. \((2,1)\) or \((3,1)\), and each adequate
partition merges \((1,2)\) with neither. The census is itself an
instance of a structure:

\begin{proposition}[adequacy as an incompatibility-graph criterion]
\label{audit-prop:adequacy}
Let \(G\) be the incompatibility graph of the target: vertices are
states, with an edge \((x, x')\) when \(R(x) \cap R(x') =
\emptyset\). A partition is adequate if and only if every fibre is an
independent set of \(G\); on the audited target the edges are exactly
\(\{(1,2)\text{--}(2,1),\ (1,2)\text{--}(3,1)\}\), so the adequate
partitions are the partitions into independent sets, counted by
inclusion--exclusion as \(15 - (5 + 5 - 2) = 7\) (the two edge events
overlap in exactly the \(2\) partitions whose fibre contains
\((1,2), (2,1), (3,1)\) jointly), and the adequate two-cell partitions
are exactly the graph's two proper \(2\)-colorings --- rows 6 and 7.
In an \(a\)-action universe, a fibre fails adequacy only through a
subfamily of at most \(a\) states with empty common intersection: the
Helly number of the family of subsets of an \(a\)-element action set
is \(a\) (verified exhaustively for \(a = 2\) and \(a = 3\), where the
minimal failing family \(\{\{1,2\}, \{2,3\}, \{1,3\}\}\) attains it);
pairwise intersection suffices only for \(a = 2\), and for convex
\(m\)-dimensional action domains the bound is \(m + 1\)
(Proposition~\ref{audit-prop:helly}).
\end{proposition}

\begin{proof}
A fibre is adequate exactly when its states' safe-action sets intersect,
which is the complement of an edge exactly on two-state fibres and,
since the audited universe has Helly number two, for all fibres
(Proposition~\ref{audit-prop:helly} is the \(m\)-dimensional analogue). The
edge count is the two disjointness facts above. The inclusion--exclusion
count: partitions whose fibres include the edge \((1,2)\text{--}(2,1)\)
are partitions of the three objects \(\{(1,2)(2,1)\}, (2,2), (3,1)\) ---
five; symmetrically five for the other edge; both edges force the
fibre \(\{(1,2), (2,1), (3,1)\}\) --- two partitions (with and without
\((2,2)\) merged). A two-cell adequate partition is a bipartition into
independent sets, i.e. a proper \(2\)-coloring of the star; there are
exactly two. The Helly statements are finite computations over all
subfamilies of the \(2^{a}\) subsets of an \(a\)-element set (exact for
\(a = 2, 3\); check P19).
\end{proof}

\subsection{Regime uncertainty}\label{audit-regimes}

Two regimes alternate: an adversary chooses at the end of step
\(t - 1\) the regime realized at step \(t\). The model, explicit: the
controller acts one instrument \(u \in \{1,2\}\) per step (coordinate
\(u\) regenerates by one, cap \(3\)), and the acting regime deals one
unit of damage to its own coordinate (floor \(0\)); there is no other
change. With the regime frozen and announced, the controller drives the state
to a fixed point --- acting the damaged coordinate holds it wherever
its level is below the cap, and one step of the other instrument frees
a capped coordinate --- so the corner \((1,1)\) is
viable under each regime separately, and each frozen kernel comprises
all \(36\) belief pairs. Mode-blind operation, with the regime read
only through the state, yields a kernel of exactly 28 --- every pair
containing the corner is lost (from \((1,1)\) either instrument
regenerates one coordinate while the adversarial damage zeroes the
other), no merged-pair loss exists, and the loss is generated entirely
at the singleton level. The mode-blind kernel
agrees setwise with the full-information kernel of Section
\ref{audit-classes}; the agreement is coincidental, arising under different
damage routings, and no identity is claimed. All three counts (frozen
\(36\) per regime, corner viable under each, mode-blind \(28\)) are
machine-verified from these primitives (check P13 of the verification script).

\subsection{The continuous benchmark}\label{audit-benchmark}

A state is a pair of stocks \((x_1, x_2) \ge 0\) with a demand floor
on each; the aggregate is \(Y = x_1 + x_2\), and the \(Y\)-fibre is
the set of states with that aggregate. One control set
\(U = \{u \ge 0 : u_1 + u_2 \ge 2\}\) faces the floors through
Y-dependent caps \(\mathrm{cap}_1(Y) = 3/2 - (Y - 2)/10\) and
\(\mathrm{cap}_2(Y) = 59/50 - (Y - 2)/10\); an allocation \(u \in U\)
is feasible on the fibre when \(u_i \le \mathrm{cap}_i(Y)\) --- fibre
feasibility, not transition dynamics, is this section's audited object.
The critical aggregate is \(Y^{*} = 27/5\): the cap sum equals \(2\)
exactly there. Table
\ref{audit-tab:bench} evaluates the caps on the audited aggregates and Figure
\ref{audit-fig:bench} displays them. The \(Y = 5\) fibre is viable with
witness \((6/5, 4/5)\); the \(Y = 6\) fibre is obstructed with cap sum
\(47/25 < 2\), certified by the dual measure \((1/2, 1/2)\) on the two
active-floor atoms with margin \(3/50\). The crossover is again a general law:

\begin{proposition}[exact benchmark crossover]\label{audit-prop:bench}
For every aggregate \(Y \ge 2\) with \(\mathrm{cap}_i(Y) \ge 0\), the
\(Y\)-fibre admits a feasible input (some \(u \in U\) with
\(u_i \le \mathrm{cap}_i(Y)\)) if and only if
\(\mathrm{cap}_1(Y) + \mathrm{cap}_2(Y) \ge 2\), which holds if and
only if \(Y \le 27/5\). For \(Y \ge 27/5\) the
\(\ell_1\)-normalized dual measure \((1/2, 1/2)\) on the two cap rows
certifies the obstruction with the exact margin
\(\tfrac12\bigl(2 - \mathrm{cap}_1(Y) - \mathrm{cap}_2(Y)\bigr) =
(Y - 27/5)/10\).
\end{proposition}

\begin{proof}
If the cap sum is at least \(2\), the input
\(u = (\mathrm{cap}_1(Y),\, 2 - \mathrm{cap}_1(Y))\) is nonnegative
(\(\mathrm{cap}_1(Y) \le 3/2 < 2\) on \(Y \ge 2\)), lies in
\(U\) (its coordinate sum is \(2\)), and respects the caps
(\(2 - \mathrm{cap}_1(Y) \le \mathrm{cap}_2(Y)\) is the cap-sum
inequality). Conversely any \(u \in U\) under the caps has
\(2 \le u_1 + u_2 \le \mathrm{cap}_1(Y) + \mathrm{cap}_2(Y)\), so a
cap sum below \(2\) is infeasible. The cap sum is the affine function
\(67/25 - (Y - 2)/5\), equal to \(2\) exactly at \(Y = 27/5\) and
declining beyond; the dual measure gives
\(\lambda^\top(\mathrm{cap} - u) = (\mathrm{cap\ sum})/2 -
(u_1 + u_2)/2 \le (\mathrm{cap\ sum})/2 - 1\), the stated margin.
\end{proof}

The derivation of the caps from a two-stock resource model, and the
fibre-width formula \(\Delta x_1(Y) = Y - 4\), are carried in the
dual-certificate companion (Abaee, 2026, Minimax dual certificates); this section owns the
per-aggregate audit and the radius reading.

\begin{table}[t]
\centering
\footnotesize
\caption{The benchmark's cap arithmetic on the audited aggregates
(exact rationals). The cap sum crosses the demand level \(2\) exactly
at \(Y^{*} = 27/5\); at \(Y = 6\) the cap sum falls \(3/25\) short of
\(2\), and the \(\ell_1\)-normalized dual measure \((1/2, 1/2)\) certifies the
obstruction with margin \(3/50\).}
\label{audit-tab:bench}
\begin{tabular}{@{}l c c c l@{}}
\toprule
\(Y\) & \(\mathrm{cap}_1\) & \(\mathrm{cap}_2\) & sum & reading\\
\midrule
4 & $13/10$ & $49/50$ & $57/25$ & above\\
$9/2$ & $5/4$ & $93/100$ & $109/50$ & above\\
5 & $6/5$ & $22/25$ & $52/25$ & viable (witness $(6/5, 4/5)$)\\
$27/5$ & $29/25$ & $21/25$ & $2$ & critical\\
$11/2$ & $23/20$ & $83/100$ & $99/50$ & below\\
6 & $11/10$ & $39/50$ & $47/25$ & obstructed, margin $3/50$\\
\bottomrule
\end{tabular}
\end{table}

\begin{figure}[t]
\centering
\includegraphics[width=0.96\linewidth]{figs_ws4/fig_benchmark.pdf}
\caption{The benchmark. Both caps decrease linearly in the aggregate
\(Y\); their sum crosses the demand level \(2\) exactly at
\(Y^{*} = 27/5\). The \(Y = 5\) fibre is viable (green), the \(Y = 6\)
fibre obstructed with margin \(3/50\) (red).}
\label{audit-fig:bench}
\end{figure}

\subsection{Static duality}\label{audit-duality}

Static duality for the benchmark's two floor rows, abstracted to drift
rows on \([0,1]\): the max--min equals the
min-\(\lambda\) max at value \(-1/10\), with multiplier
\(\lambda = (1/2, 1/2)\) the \(\ell_1\)-normalized Farkas pair of the calculus's
worked certificate (Abaee, 2026, An obstruction calculus). Figure \ref{audit-fig:duality} displays
the instance: under the \(\ell_1\)-normalized multiplier the expected drift is
constant \(-1/10\) on all of \([0,1]\). Convexity is load-bearing: with
\(U = \{0, 1\}\) and rows \((-1, +1)\), \((+1, -1)\), the common set is
empty, the max--min is \(-1\), and the min-multiplier maximum is
\(0\). Under convexity the witness is sparse:

\begin{proposition}[static minimax duality on the two-floor instance]
\label{audit-prop:saddle}
On the obstructed instance with drift rows \(\psi_1(u) = 2/5 - u\),
\(\psi_2(u) = u - 3/5\) on \(U = [0,1]\),
\[
\max_{u \in [0,1]} \min_i \psi_i(u) \;=\;
\min_{\lambda \in \Delta_2} \max_{u \in [0,1]}
\bigl(\lambda_1 \psi_1 + \lambda_2 \psi_2\bigr) \;=\; -\,\frac{1}{10},
\]
with the unique saddle point \((u^{*}, \lambda^{*}) = \bigl(1/2,
(1/2, 1/2)\bigr)\); under \(\lambda^{*}\) the expected drift is
constant \(-1/10\) on all of \([0,1]\).
\end{proposition}

\begin{proof}
The min envelope \(\min_i \psi_i(u)\) is concave and piecewise
linear with its peak at the crossing \(2/5 - u = u - 3/5\), i.e.\
\(u^{*} = 1/2\), value \(-1/10\). For \(\lambda \in \Delta_2\)
the pooled drift \((2\lambda_1 - 1)\cdot(-u) + \text{const}\) is
affine with slope \(\lambda_2 - \lambda_1 = 1 - 2\lambda_1\); it is
constant exactly at \(\lambda^{*} = (1/2, 1/2)\), where the value is
\(\tfrac12(2/5) + \tfrac12(-3/5) = -1/10\) for every \(u\). At any
other \(\lambda\) the maximum over \([0,1]\) is attained at an
endpoint; writing \(A(\lambda)\) for the \(u = 0\) endpoint value, the
two endpoint values are \(A(\lambda)\) and \(A(1 - \lambda)\), and
the averaging bound
\(\max(A(\lambda), A(1-\lambda)) \ge \tfrac12\bigl(A(\lambda) +
A(1-\lambda)\bigr) = -\tfrac1{10}\) --- the average is exactly the
constant value, for every \(\lambda\) --- with equality only at
\(\lambda^{*}\): the convexity of the max, not symmetry alone, locates
the minimum. Weak duality (max-min \(\le\) min-max)
holds unconditionally; the exhibited pair attains both sides at
\(-1/10\), which is therefore the exact common value and the unique
saddle.
\end{proof}

\begin{proposition}[sparse obstruction witnesses]\label{audit-prop:helly}
An infeasible finite family of affine constraints on a convex domain in
\(\mathbb{R}^m\) has an infeasible subfamily of at most \(m + 1\)
members, and \(m + 1\) is tight. The instance realizes the bound: the
three constraints \(\{u_1 \ge 1,\ u_2 \ge 1,\ u_1 + u_2 \le 1\}\)
are pairwise feasible, jointly infeasible, and no two of them certify
the obstruction.
\end{proposition}

\begin{proof}
Intersect each constraint's halfspace with the domain: an infeasible
family of compact convex sets in \(\mathbb{R}^m\) has an infeasible
subfamily of size at most \(m + 1\) (Helly's theorem), and \(m + 1\)
is tight: the reversed halfspaces \(\{u_i \ge 1\}\) and
\(\sum_i u_i \le 1\) in \(\mathbb{R}^m\) --- the facets of an
\(m\)-simplex, translated outward --- meet \(m\) at a time and never
all \(m+1\) (the facets themselves meet in the simplex, which is why
the reversal is needed). For the
instance, each pair is feasible --- at \((1,1)\), \((1,0)\), and
\((0,1)\) respectively --- while all three give the contradiction
\(2 \le 1\); the verification script verifies each pair at its printed
feasible point and the triple by the exact contradiction.
\end{proof}

The Helly-tight family of three constraints
\(\{u_1 \ge 1,\ u_2 \ge 1,\ u_1 + u_2 \le 1\}\) --- no two of them
already infeasible --- attains the obstruction with \(m{+}1 = 3\)
constraints. The feasible contrast at rows
\((2/5 - u,\ u - 1/5)\) carries value \(+1/10\) on both sides of the
saddle: a certificate exists exactly where the value is negative, and at
a zero value the instance is boundary-critical with no strictly positive
margin available.
Weak duality (max--min \(\le\) min--max) holds unconditionally. Strong duality holds where its hypotheses are verified: the control set compact \emph{and convex}, the rows affine (or concave) in \(u\), and the label family finite --- as in the static instance above and the finite-horizon static-observation classes of the calculus (Abaee, 2026, An obstruction calculus). Both hypotheses are load-bearing: the paper's own nonconvex example (\(U = \{0,1\}\), rows \((-1,+1)\), \((+1,-1)\)) is compact, finite, and row-affine, yet gaps (\(-1\) against \(0\)); and continuity of the rows does not repair convexity, as the nonconvex rows example above shows --- concavity of the rows in \(u\), not continuity, is what the minimax argument consumes.

\begin{figure}[t]
\centering
\includegraphics[width=0.96\linewidth]{figs_ws4/fig_duality.pdf}
\caption{Static duality on the two-floor instance. The two drift rows
cross at \(u = 1/2\); the max--min equals the min-multiplier maximum at
\(-1/10\), and the \(\ell_1\)-normalized Farkas pair \((1/2, 1/2)\) makes the
expected drift constant at the value.}
\label{audit-fig:duality}
\end{figure}

\subsection{The certainty-equivalence drift audit}\label{audit-drift}

For the index \(g(s) = s^{2}\) the uncorrected law's one-step drift is
\(g(s + 1/10) - g(s) = s/5 + 1/100\) on \([1,2]\): increasing, with
minimum \(21/100\) at \(s = 1\) and maximum \(41/100\) at \(s = 2\)
(Table \ref{audit-tab:drift}); the corrected law's drift is identically zero --- by
construction, the corrected law being the lead removed. A biased
observation map misreads the state by a fixed one-step lead; on any
decision model conditioned on the map, that misreading is the mechanism
through which a nonempty kernel could empty, and the drift audit bounds
it exactly --- the remedy is structural correction of the map, not a
better sensor.

\begin{table}[t]
\centering
\footnotesize
\caption{The certainty-equivalence drift audit:
\(g(s + 1/10) - g(s) = s/5 + 1/100\) exactly; the corrected law's
drift is \(0\) at every probed state.}
\label{audit-tab:drift}
\begin{tabular}{@{}l c c c@{}}
\toprule
\(s\) & uncorrected & corrected & increment per \(1/4\)\\
\midrule
1 & $21/100$ & 0 & \\
$5/4$ & $26/100$ & 0 & $1/20$\\
$3/2$ & $31/100$ & 0 & $1/20$\\
$7/4$ & $36/100$ & 0 & $1/20$\\
2 & $41/100$ & 0 & $1/20$\\
\bottomrule
\end{tabular}
\end{table}

\subsection{Multiple floors}\label{audit-floors}

With several binding floors the Farkas certificate obstructs jointly: a
single certificate over all floors certifies nonviability where
per-floor certificates do not compose. On the audited benchmark's two
floors (Section \ref{audit-benchmark}), the minimal infeasible subsystem
carries explicit witnesses --- the Helly-tight triple of Section
\ref{audit-duality} --- and the corner \((1,1)\) is dynamically nonviable
as in the shared system. \emph{Quantifier and dual form.} Joint maintenance of several floors is the quantifier \(\exists u\ \forall j \in I(x)\) --- one control satisfying every active floor --- not \(\forall j\ \exists u_j\); per-floor certificates do not compose. The gap is visible in dimension one: on \(U = [-1, 1]\) the floors \(g_1(u) = u - \tfrac12 \ge 0\) and \(g_2(u) = -u - \tfrac12 \ge 0\) are each separately maintainable, yet no common control maintains both, and the \(\ell_1\)-normalized multiplier \(\lambda = (\tfrac12, \tfrac12)\) certifies the joint infeasibility with \(\sup_{u \in U}(\lambda_1 g_1 + \lambda_2 g_2)(u) = -\tfrac12 < 0\) --- a strict-separation dual certificate of exactly the Proposition~\ref{audit-prop:helly} shape, with \(m + 1 = 2\) constraints necessary in dimension \(m = 1\) (each single floor is feasible). Scalar aggregation inherits the gap: a positive weighted average \(\sum_j \lambda_j h_j\) can conceal a violated component, and a multiplier inequality \(\sum_j \lambda_j \mathrm{D}h_j \ge 0\) does not imply that every active drift is nonnegative; the \(\ell_1\)-normalized dual measure of Section \ref{audit-benchmark} certifies the aggregate shortfall exactly because both cap rows are charged jointly. In the audited convex static examples, infeasibility is certified by a multiplier supported on at most \(m + 1\) active constraints; no such bound is claimed for nonconvex controls or dynamic-policy images. Two further objects are recorded for subsequent study: a
first-breach timing bound for multi-floor systems, and a corner-cone
reading of the infeasibility region.

\subsection{Robustness margins}\label{audit-margins}

Each certificate carries a margin, and each margin is a robustness
radius --- the size of the perturbation the certificate's verdict
tolerates.

\begin{proposition}[margins as radii]\label{audit-prop:margins}
(i) \emph{Benchmark}: at \(Y = 6\) the obstruction survives any
per-cap increase \(\delta < 3/50\) of both caps and only there: the
certificate's margin \(3/50\) is exactly the per-cap robustness radius
(feasibility returns at \(\delta = 3/50\)). (ii) \emph{Timing}: the
boundary cells have zero slack (\(z_0 - 1 - T_{\mathrm{obs}} = 0\)):
any adversarial damage above unit rate flips them --- the inclusive
convention is exactly the statement that the audited identity holds on
zero-margin cells. (iii) \emph{Drift}: the corrected law's zero drift
has the sign of the correction error \(\varepsilon\) (\(2\varepsilon +
\varepsilon^2 = \varepsilon(2 + \varepsilon)\) at \(s = 1\)): the
sign-certificate is robust for every \(\varepsilon \in (-2, 0) \cup
(0, \infty)\). (iv) \emph{Decentralization}: each of the four forced
votes of the audited belief is knife-edge --- flipping that single vote
sustains nothing (zero margin, exact); the protocol's value is carried
by four coordinates, each of which is load-bearing.
\end{proposition}

\begin{proof}
(i) feasibility at \(Y = 6\) under uniform increases \(\delta\) is
\(\tfrac{47}{25} + 2\delta \ge 2 \iff \delta \ge
\tfrac{3}{50}\) --- exact arithmetic, checked at \(\delta =
\tfrac{3}{200}\) (infeasible) and \(\delta = \tfrac{3}{50}\)
(feasible). (ii) at \((z_0, T) = (2, 1)\) the worst damage under rate
\(1 + \varepsilon\) is \(1 + \varepsilon > 1 = z_0 - 1\). (iii) the
sign evaluation is the displayed factorization. (iv) for each forced
cell, the count of sustaining law pairs with that vote flipped is zero
(exhaustive over the \(256\) pairs; check P22).
\end{proof}

\subsection{Design rules}\label{audit-rules}

The rules hold on the audited systems --- each an exact census identity or a proved general law, tagged by its evidence --- and transfer as default guidance only at that evidence level. \emph{Timing:} (proved general law, Proposition \ref{audit-prop:timing}) reviews are boundary-inclusive ---
\(T_{\mathrm{obs}} \le \tau^{*}(z_0)\), the deadline at worst-case exit admissible under the audited endpoint convention (Section \ref{audit-timing}). \emph{Coarseness:} (exact census, Section \ref{audit-adequacy}) a fibre
merging states with disjoint safe-action sets obstructs at that
information state; the remedy is a fibre split, not a finer threshold
on the same index (Section \ref{audit-adequacy}). \emph{Aggregation:} (exact census, Sections \ref{audit-classes} and \ref{audit-adequacy}) the aggregate index's fibres cross the safe boundary --- each mixes states of conflicting verdict --- and the aggregate kernel (\(26\) against the full \(28\)) certifies only a certainly-safe inner approximation; the certainly-safe set is the sound report.
\emph{Bias:} (exact audit, Section \ref{audit-drift}) the uncorrected law's drift is at least \(21/100\) on a
101-point grid of \([1,2]\), while the corrected law's drift is
identically zero; a fixed bias is exactly correctable, and the remedy is
structural correction of the map, not a better sensor. The decentralized audit adds a final rule: (proved --- the restricted kernel is zero, Proposition \ref{audit-prop:authority}) a whole-law authority restriction on one agent's vote collapses the sustainable set outright, so protocol
design precedes instrument tuning (Section \ref{audit-decentralized}).

\subsection{Mechanized verification}
\label{audit-mech}

The master monotonicity statement of Proposition~\ref{audit-prop:master} is formalized in the
project's Lean~4 layer, in namespace \texttt{Formalizations.WS}
(\texttt{WS\_WorkedSystems}): \texttt{master\_monotone} is the theorem itself, and
\texttt{monotone\_action\_axis} the policy axis of it, both stated over the abstract audited
tuple rather than over any particular audited instance, so they hold for every system in the
family at once.

\noindent\textbf{No admitted gaps.} A comment-stripped source scan of the whole project finds
no \texttt{sorry}, \texttt{admit}, \texttt{axiom} or \texttt{constant} declaration anywhere.

\noindent\textbf{Scope, stated plainly.} What is machine-checked is the monotonicity
theorem --- the one abstract statement whose instances the audits tabulate. \textbf{The audits
themselves are not formalized}: the counts of survivable belief pairs, the kernel sizes under
policy-class restrictions, the partition-adequacy census and the dual certificates are
computations performed by the deposited scripts in integer and rational arithmetic, and are
regenerated rather than proved. This is the mirror image of the usual arrangement, and it is
worth being exact about it: the theorem that organises the table is proved, while the entries
in the table are computed.

\subsection{Verification methods}\label{audit-methods}

Every count, witness, identity, bound, figure, and tabulated entry in
this paper is regenerated by deterministic scripts in exact integer and
rational arithmetic, standard library only. The master verification
script recomputes the counts of Sections \ref{audit-master}--\ref{audit-regimes}
from the shared system's primitives --- the transition, the fibre
partition function, and the backward recursion with the no-repeat
protocol threaded as a parameter --- re-derives every tabulated entry of
Tables \ref{audit-tab:master}, \ref{audit-tab:laws}, \ref{audit-tab:timing},
\ref{audit-tab:census}, \ref{audit-tab:bench}, and \ref{audit-tab:drift} independently,
enumerates the census of Section \ref{audit-adequacy}, evaluates the
benchmark identities of Section \ref{audit-benchmark} in rational
arithmetic, asserts the drift bounds, and checks that each headline
number appears in the text. The figure-and-table generation script
recomputes and asserts each plotted or printed quantity from the same
primitives before drawing; the two scripts agree on every shared
quantity. The audited identities are additionally covered by the
per-system verification scripts of the companion papers (Abaee, 2026, An obstruction calculus;
see the Code availability statement). Two independent recomputation
paths guard the central recursions (checks P17--P18): an exhaustive
enumeration of all deterministic stationary policies --- \(2^{5} = 32\)
aggregate, \(2^{9} = 512\) full-information --- whose union of
safe-pair sets reproduces the kernels exactly (memoryless determinacy
on the instance, itself an exact verified fact), and a fixpoint
computation of each class's horizon family, whose stabilization equals
the audited kernel setwise. The recursion's cost is measured exactly:
the memoized backward recursion over the \(36\)-pair population uses
the distinct-subproblem counts printed by the script (states \(\times\)
horizon \(\times\) register per class), the policy enumerations use
\(32\) and \(512\) candidates, the decentralized census \(256\) law
pairs (\(64\) behavioral), the adequacy census \(15\) partitions, and
the two timing grids \(93\) cells each.

\subsection{Conclusion}\label{audit-conclusion}

Exact audits bind the obstruction calculus to its worked systems: each
mechanism acquires a count, a witness, and a boundary. The timing
certificate fires everywhere beyond \(z_0 - 1\) and nowhere at the
boundary; the adequate monitoring designs on the audited target are
enumerable, and no coarsest one exists; policy-class value drops
unequally across aggregation and protocol restrictions, with the two
middle cells incomparable, the meet of their kernels institutional
setwise, and exactly one pair --- the alternation-forced pair --- beyond
both; unanimity costs a fixed and identified set of beliefs, sustained
by exactly the \(4 \times 4\) product of free votes, recoverable only
while both votes remain free; regime blindness loses exactly the
corner's pairs, while no frozen regime loses any; and the continuous
benchmark separates the viable from
the obstructed at \(Y^{*} = 27/5\) with rational certificates on both
sides. The incomparability of the middle cells is a shared-register
phenomenon: with per-branch registers the protocol cost collapses
(Proposition~\ref{audit-prop:register}), which is precisely why the shared
register is the audited institutional object. The design rules --- boundary-inclusive reviews, fibre splits
over threshold tuning, certainly-safe reporting under aggregation,
structural repair of bias, and protocol design before instrument tuning
--- are the audit's operational content, and the master table is its
evidence: every count in the paper is one column of one table, every entry of which is regenerated by script. These examples do not constitute a general taxonomy of viability failures; they supply exact counterexamples, finite censuses, and reproducible benchmarks that separate mechanisms often conflated in qualitative discussion. Read through the master theorem, the audit is a small theory: the five mechanisms --- common-action failure, memory/register coupling, finite-time boundary exit, decentralized information loss, and joint convex infeasibility --- each carry a certificate, a monotonicity status under the theorem's three axes, a requirement-level characterization (statewise factoring for unanimity, incompatibility-graph adequacy for fibres, max-plus hitting times for regimes), and a robustness radius (\(3/50\) per cap on the benchmark, zero on the boundary cells and the forced votes, sign-robustness on the drift). Refinement completeness on the continuous side is the matching axis: the certification companion's meshes, stored scenarios, and belief cells tighten its sandwich exactly as the three axes here enlarge the kernels, with positivity non-monotone on both sides (Abaee, 2026, Computational viability certification)

\section{Cross-part conclusion}
\label{conclusion}

The two parts above make one argument between them: \emph{exactness is the instrument, not the
limitation.} That claim has two independent supports, and it is worth separating them, because they
fail independently and are established independently.

\emph{Exactness scales.} Part I takes the discipline into the regime where it is expected to break
and reports what happens. The reason it does not break is structural rather than computational: the
stored object is an \emph{antichain} and the argued instrument is a \emph{pairwise bound}, and
between them $1{,}048{,}576$ subset evaluations collapse to $496$ sets without loss --- and, at five
parameters, $68{,}719{,}476{,}736$ evaluations onto $1{,}552$ sets. This is not a constant-factor
saving; it is a change of what is being enumerated. The combinatorial structure is doing the work:
a pairwise-Hamming lemma forces every blind survivable set to consist of Hamming-adjacent cells,
and no three cells of the four-cube are pairwise adjacent, so \emph{no blind policy of any period
keeps three or more cells safe}. The classification is therefore exact and complete rather than
sampled. Part I also reports where the instrument runs out --- at five parameters the pairwise bound
\emph{stops being sufficient} --- which is the honest form of a scaling claim.

\emph{Exactness bites.} Part II shows the other half: the verdicts that exactness produces are not
the ones a looser treatment would have predicted, and several of them are negative in a way that
matters for design. The review-timing identity is \emph{non-strict}: viability at the deadline holds
if and only if $T_{\mathrm{obs}} \le \tau^{*}(z_0) = z_0 - 1$, boundary included, while the strict
form against the real exit bound fails exactly on the boundary cells. On the monitoring target, the
adequate partitions are enumerable but \emph{no coarsest adequate partition exists} --- so there is
no canonical ``cheapest adequate design'' to search for, and the search for one is ill-posed. Two
middle cells are simply incomparable, and that incomparability is a \emph{shared-register}
phenomenon: with per-branch registers the protocol cost collapses, which is precisely why the shared
register is the audited institutional object. Each of these is a statement about a boundary, a
non-existence, or an incomparability --- and each is exactly the kind of statement that a
floating-point treatment would smooth into a wrong answer.

\emph{Why the two belong together.} A method that scales is not thereby shown to produce
surprising verdicts, and a collection of surprising exact examples is not thereby shown to survive
dimensionality. Established together, they close the two objections that are normally raised
against exact methods: that they are confined to toys, and that they are a stylistic preference
rather than a source of results. Part I answers the first; Part II answers the second. The shared
commitment that makes both possible is procedural as much as mathematical --- every count in both
parts is regenerated by script, over integer and rational arithmetic, so each number in the tables
is open to re-inspection rather than taken on trust. \emph{That is what turns a computation into
evidence.}

\emph{What is deliberately not claimed.} These examples do not constitute a general taxonomy of
viability failures. They supply exact counterexamples, finite censuses and reproducible benchmarks
that separate mechanisms often conflated in qualitative discussion. Read through the master theorem
of Part II, the audit is a small theory: five mechanisms --- common-action failure,
memory/register coupling, finite-time boundary exit, decentralized information loss, and joint
convex infeasibility --- each carry a certificate, a monotonicity status under three axes, and a
requirement-level characterization. But the taxonomy is not exhaustive, and neither part claims
otherwise.

\emph{The operational content.} Part II draws the design rules the audit licenses, and they are
worth restating because they are the transferable output: boundary-inclusive reviews; fibre splits
over threshold tuning; certainly-safe reporting under aggregation; structural repair of bias rather
than compensation for it; and protocol design before instrument tuning. The last of these is the
general lesson both parts support --- \emph{design the observation before tuning the estimator},
because in both the scaling regime and the audit regime it is the information structure, not the
estimation quality, that determines what is achievable.

\subsection*{References}
\label{references}
Abaee, A. (2026).. An obstruction calculus for viability under incomplete observation. Submitted.

Abaee, A. (2026). Computational viability certification under incomplete
observation: a continuous-to-finite bridge with exact certificates, a
solver-assisted certified campaign, and a finite-side reference library.
Submitted.

Abaee, A., 2026. Minimax dual certificates for the common-action
obstruction. Manuscript submitted for publication.

Abaee, A. (2026).. Probabilistic sufficiency for the obstruction calculus: belief-state safety values under partial observation. Submitted.

Abaee, A., 2026. Robust viability of the 2J3KL limit reference point
under a surplus-production map: policy scoring, expansion, and when
catch cannot help. Zenodo.
\url{https://doi.org/10.5281/zenodo.22552060}

Abaee, A., 2026. An obstruction calculus for viability under
incomplete observation (submitted).

Althoff, M., et al., 2020. ARCH-COMP20 category report: continuous and hybrid systems with
linear continuous dynamics. In: \emph{ARCH20, 7th International Workshop on Applied
Verification of Continuous and Hybrid Systems}, EPiC Series in Computing, vol.~74, pp.~16--48.

Aubin, J.-P., 1991. \emph{Viability Theory}.

Baccelli, F., Cohen, G., Olsder, G.J. and Quadrat, J.-P., 1992. \emph{Synchronization and
Linearity: An Algebra for Discrete Event Systems}.

Baccelli, F., Cohen, G., Olsder, G.J., Quadrat, J.-P., 1992.
Synchronization and Linearity: An Algebra for Discrete Event Systems.

Birkh\"auser, Boston.

Chapter 4, The Sperner
property. \url{https://math.mit.edu/~rstan/algcomb/algcomb.pdf}.

Engel, K., 1997. \emph{Sperner Theory}. Cambridge University Press.

%% ------------------------------------------------------------------
%% Entries added 2026-09-29 to support subsection \ref{priorart}.
%% Chatterjee/Doyen/Henzinger copied verbatim from the sibling paper's
%% bibliography (already cited there), so it is consistent across the set.
%% NEEDING PUBLISHER VERIFICATION BEFORE SUBMISSION (pagination or exact
%% venue not confirmed here):
%%   Sperner 1927; Kleitman 1969; Pineau et al. 2003 (pages);
%%   Shani et al. 2013 (pages).
%% ------------------------------------------------------------------
Chatterjee, K., Doyen, L., Henzinger, T.A., 2009. Qualitative analysis of
partially-observable Markov decision processes. In: Int. Symp. Math. Found. Comput. Sci.
(MFCS 2009).

Farkas, J., 1902. Theorie der einfachen Ungleichungen. \emph{Journal f\"ur die Reine und
Angewandte Mathematik} \textbf{124}, 1--27.

Geretti, L., Dit Sandretto, J.A., Althoff, M., et al., 2020. ARCH-COMP20 category report:
continuous and hybrid systems with nonlinear dynamics. In: \emph{ARCH20, 7th International
Workshop on Applied Verification of Continuous and Hybrid Systems}, EPiC Series in Computing,
vol.~74, pp.~49--75.

Helly, E., 1923. \"Uber Mengen konvexer K\"orper mit gemeinschaftlichen Punkten.
\emph{Jahresbericht der Deutschen Mathematiker-Vereinigung} \textbf{32}, 175--176.

Kleitman, D.J., 1969. On Dedekind's problem: the number of monotone Boolean functions.
\emph{Proc. Amer. Math. Soc.} \textbf{21}, 677--682.

LNCS, vol. 5734, pp. 635--646.

Lovejoy, W.S., 1991. Computationally feasible bounds for partially
observed Markov decision processes.

Minato, S. (1993). Zero-suppressed BDDs for set manipulation in combinatorial problems. In: Proceedings of the 30th Design Automation Conference (DAC '93), 272--277.

Operations Research 39, 162--175.

Pineau, J., Gordon, G., Thrun, S., 2003. Point-based value iteration: an anytime algorithm
for POMDPs. \emph{Proc. IJCAI 2003}.

Shani, G., Pineau, J., Kaplow, R., 2013. A survey of point-based POMDP solvers.
\emph{Autonomous Agents and Multi-Agent Systems} \textbf{27}(1), 1--51.

Sperner, E., 1927. Ein Satz \"uber Untermengen einer endlichen Menge. \emph{Math. Z.}
\textbf{27}, 544--548.

Stanley, R.P., 2013. \emph{Topics in Algebraic Combinatorics}.

Wiley, Chichester.

\subsection*{Declarations}

\subsection*{Funding}
None.

\subsection*{Competing interests}
None.

\subsection*{Data availability}
No external data were used.

\subsection*{Code availability}
The verification script
\url{paper2_exact_belief_computation_v13_verification.py} is publicly
available at \url{https://github.com/MIKEAA2020/general-sustainability}
(folder \texttt{arena agent 1/paper rewrites/latex}); it reproduces all bounds, bands, ladder values, pairings,
census counts, and thresholds verbatim.

\subsection*{AI declaration}
GLM (Z.ai), Qwen (Alibaba Cloud) and DeepSeek AI assisted with drafting
and iterative review. The author reviewed and edited outputs and takes
responsibility for the final work.
\subsection*{Declarations}

\subsection*{Funding}
None.

\subsection*{Competing interests}
None.

\subsection*{Data availability}
No external data were used.

\subsection*{Code availability}
The verification script \url{paper2_worked_systems_v17_verification.py}
and the figure-and-table generation script
\url{paper2_worked_systems_figures.py} are publicly available at
\url{https://github.com/MIKEAA2020/general-sustainability} (folder
\texttt{arena agent 1/paper rewrites/latex}); they reproduce all
counts, witnesses, bounds, figures, and tabulated entries verbatim.

\subsection*{AI declaration}
GLM (Z.ai), Qwen (Alibaba Cloud) and DeepSeek AI assisted with drafting
and iterative review. The author reviewed and edited outputs and takes
responsibility for the final work.

\end{document}
